\documentclass[11pt,a4paper]{article}
\usepackage{fontspec}
\setmainfont[BoldFont={Noto Sans CJK JP Bold}]{Noto Serif CJK JP}
\setsansfont[BoldFont={Noto Sans CJK JP Bold}]{Noto Sans CJK JP}
\usepackage{amsmath,amssymb,amsthm,mathtools}
\usepackage[margin=24mm]{geometry}
\usepackage{xcolor,booktabs,longtable,array,enumitem,fancyhdr,titlesec}
\usepackage[colorlinks=true,linkcolor=blue!50!black,citecolor=blue!50!black,urlcolor=blue!50!black]{hyperref}
\setlength{\parindent}{1em}
\setlength{\parskip}{2pt}
\linespread{1.17}
\emergencystretch=4em
\setlist{itemsep=1pt,topsep=3pt}
\setlength{\tabcolsep}{4pt}

%%% --- Japanese line breaking for XeTeX (no xeCJK available) ---
\makeatletter
\XeTeXinterchartokenstate=1
\newXeTeXintercharclass\cjkc
\newXeTeXintercharclass\cjkclose
\newXeTeXintercharclass\cjkopen
\def\setcjkrange#1#2{\@tempcnta=#1 \@whilenum\@tempcnta<#2 \do{\XeTeXcharclass\@tempcnta=\cjkc \advance\@tempcnta by 1 }}
\setcjkrange{"3000}{"30FF}
\setcjkrange{"4E00}{"9FFF}
\setcjkrange{"FF00}{"FFEF}
\setcjkrange{"2010}{"2027}
\setcjkrange{"25A0}{"25FF}
\setcjkrange{"2190}{"21FF}
\@tempcnta=0
\def\setclose#1{\XeTeXcharclass#1=\cjkclose}
\def\setopen#1{\XeTeXcharclass#1=\cjkopen}
\setclose{"3001}\setclose{"3002}\setclose{"FF0C}\setclose{"FF0E}\setclose{"FF09}\setclose{"300D}\setclose{"300F}\setclose{"3011}\setclose{"FF1A}\setclose{"FF1B}\setclose{"FF01}\setclose{"FF1F}\setclose{"30FC}
\setopen{"FF08}\setopen{"300C}\setopen{"300E}\setopen{"3010}
\def\cjkglue{\hskip0pt plus 0.25pt minus 0.1pt\relax}
\XeTeXinterchartoks\cjkc\cjkc={\cjkglue}
\XeTeXinterchartoks\cjkclose\cjkc={\cjkglue}
\XeTeXinterchartoks\cjkclose\cjkopen={\cjkglue}
\XeTeXinterchartoks\cjkc\cjkopen={\cjkglue}
\XeTeXinterchartoks\cjkclose\cjkclose={}
\XeTeXinterchartoks\cjkc\cjkclose={}
\XeTeXinterchartoks 0 \cjkc={\cjkglue}
\XeTeXinterchartoks \cjkc 0={\cjkglue}
\XeTeXinterchartoks 0 \cjkopen={\cjkglue}
\XeTeXinterchartoks \cjkclose 0={\cjkglue}
\makeatother

%%% --- status tags ---
\newcommand{\tP}{\textcolor{green!45!black}{\textbf{【証明済】}}}
\newcommand{\tC}{\textcolor{blue!55!black}{\textbf{【文献】}}}
\newcommand{\tV}{\textcolor{violet!80!black}{\textbf{【計算で確認】}}}
\newcommand{\tK}{\textcolor{red!70!black}{\textbf{【予想】}}}
\newcommand{\tQ}{\textcolor{orange!85!black}{\textbf{【要確認】}}}

%%% --- theorem-like environments ---
\theoremstyle{definition}
\newtheorem{defn}{定義}[section]
\newtheorem{thm}[defn]{定理}
\newtheorem{prop}[defn]{命題}
\newtheorem{lem}[defn]{補題}
\newtheorem{cor}[defn]{系}
\newtheorem{ex}[defn]{例}
\newtheorem{rem}[defn]{注意}
\newtheorem{obs}[defn]{観察}
\newtheorem{conj}[defn]{予想}
\newcommand{\src}[1]{\par\noindent{\small\textcolor{blue!55!black}{\textbf{[出典]}~#1}}\par}
\renewcommand{\proofname}{証明}

\newenvironment{hyptab}{\par\medskip\noindent{\small\textbf{仮定と、その役割}}\par\nopagebreak\small\noindent\begin{tabular}{>{\raggedright\arraybackslash}p{3.3cm}>{\raggedright\arraybackslash}p{5.5cm}>{\raggedright\arraybackslash}p{5.9cm}}\toprule 仮定 & なぜ必要か & 証明のどこで効くか\\\midrule}{\\\bottomrule\end{tabular}\par\medskip}

%%% --- notation ---
\newcommand{\R}{\mathbb R}\newcommand{\C}{\mathbb C}\newcommand{\Z}{\mathbb Z}
\newcommand{\bT}{\mathbb T}
\newcommand{\Lb}{\overline L}
\newcommand{\Jc}{\mathcal J}
\newcommand{\dbar}{\bar\partial}
\newcommand{\ch}{\mathrm{ch}}\newcommand{\Td}{\mathrm{Td}}
\newcommand{\rk}{\mathrm{rk}}
\newcommand{\HH}{\mathrm{HH}}
\newcommand{\Res}{\mathrm{Res}}
\newcommand{\dlog}{\mathrm{dlog}\,}
\newcommand{\sh}{\sharp}

\pagestyle{fancy}\fancyhf{}
\fancyhead[L]{\small 一般化複素幾何の「三種の神器」：具体例からのアプローチ}
\fancyhead[R]{\small 調査ノート 2026-11-22}
\fancyfoot[C]{\small\thepage}
\titleformat{\section}{\large\bfseries\sffamily}{\thesection}{0.8em}{}
\titleformat{\subsection}{\normalsize\bfseries\sffamily}{\thesubsection}{0.8em}{}

\begin{document}
\sloppy

\begin{center}
{\LARGE\bfseries\sffamily 正則 Poisson 構造に付随する一般化複素構造上の\\[3pt]一般化正則直線束のコホモロジー}\\[8pt]
{\large Serre 型消滅・双対性・Riemann--Roch 型公式への\\具体例からのアプローチ}\\[12pt]
{\large Claude（調査・執筆）\ ／\ 森（問題設定）}\\[3pt]
2026年11月22日（\textbf{v2}）\\[3pt]
{\small （個人用の調査ノート。査読を受けておらず、プレプリントでもありません。\\
証明の細部と引用の書誌情報には未確認の箇所があり、本文中に\tQ\ で明示してあります。）}
\end{center}

\begin{abstract}
\noindent
複素幾何の「三種の神器」（Hirzebruch--Riemann--Roch の定理、Serre の消滅定理、Kodaira の埋め込み定理）の一般化複素幾何版を、最も簡単な非自明な例から攻めるという方針で調べた結果をまとめる。
対象は、正則 Poisson 構造 $\sigma$ から $\beta$-変換で得られる一般化複素構造 $\Jc_\sigma$（$\sigma$ が消える因子 $D$ の上で型が跳ぶ）と、その上の一般化正則直線束 $L$ の一般化 Dolbeault コホモロジー $\HH^*(L)$ である。
主な結果は次の通り。
(1) $\sigma$ が $D$ に沿って対数シンプレクティックのとき、$L$ の一般化正則構造は $D$ に沿った対数的平坦接続と同じものであり、留数 $r_i$ は定数で $c_1(L)=-\sum r_i[D_i]$ となる。
(2) 曲面で $D$ が通常二重点のみをもつ場合、$\HH^k(L)=H^k(X;\mathcal F)\oplus\C^{\mu}[-2]$（$\mathcal F$ は留数で決まる構成可能複体、$\mu$ は二重点の個数）。全 Euler 標数は $\chi(X)$ で $L$ によらず、$L$ への依存は Poincaré 多項式のみに現れる。
(3) 任意次元の射影多様体上の $\beta$-変換構造では、ample な Poisson 加群 $L_0$ について $m\gg0$ のとき横断 Dolbeault 次数 $a>0$ の項が消え、$\HH^k(L_0^m)=0\ (k>n)$ となる（Serre 型定理）。全次数での消滅は $\sigma=0$ でも成り立たない。
(4) 任意のコンパクト複素多様体・任意の正則 Poisson 構造で $\HH^k(L)^*\cong\HH^{2n-k}(L^\vee\otimes K_X^2)$ が成り立つ。
(5) トーリックな場合の weight 公式、4 次元での $\sigma$ の算術への依存、三重点での局所項の計算（一部は予想）。
一般化 Kodaira 埋め込み定理および一般の一般化 Riemann--Roch 公式は得られていない。指数理論的な見通し（b-計算、指標根）は予想の段階である。
\end{abstract}

\paragraph{本ノートの読み方。}
各記述に次のタグを付ける。\tC\ 文献からの引用（出典を明記）、\tP\ 筆者が証明したもの（標準的事実を引用した上で、主要なものは証明または証明の概略を載せる）、\tV\ 厳密な線形代数などの計算機による確認（証明ではない）、\tK\ 予想・見通し、\tQ\ 未確認の点（書誌情報を記憶に頼った箇所、引用の適用範囲が未確認の箇所など）。
\tP\ と付けた命題でも、前提として引用した文献の結果に依存している。依存関係は \S\ref{sec:ledger} にまとめた。

\tableofcontents

%--------------------------------------------------------------
\section{記号表}\label{sec:notation}
%--------------------------------------------------------------
一般的すぎる記号（外積、微分など）は省く。

\begin{longtable}{>{\raggedright\arraybackslash}p{3.4cm}>{\raggedright\arraybackslash}p{11.6cm}}
\toprule 記号 & 意味\\ \midrule\endhead
$X$ & コンパクト複素多様体、$\dim_\C X=n$（対数シンプレクティックの場合は $n$ は偶数）\\
$T,\ T^*,\ \Omega^p$ & 正則接束、正則余接束、正則 $p$-形式の層。$K_X=\Omega^n$\\
$\bT=T\oplus T^*$ & 一般化接束。内積 $\langle X+\xi,Y+\eta\rangle=\tfrac12(\xi(Y)+\eta(X))$、Courant（Dorfman）括弧 $[X+\xi,Y+\eta]=[X,Y]+\mathcal L_X\eta-\iota_Yd\xi$\\
$\Jc$ & 一般化複素構造。$\Jc^2=-1$、内積を保ち、$\pm i$ 固有束が Courant 括弧で閉じる\\
$L$ & $\Jc$ の固有束（$\bT_\C$ の極大迷零部分束、複素 Lie 亜群）。$L^*\cong\Lb$\\
$\sigma$ & 正則 Poisson 構造：$\sigma\in H^0(X,\wedge^2T)$、$[\sigma,\sigma]=0$（Schouten 括弧）\\
$\sigma^\sh$ & $T^*\to T$、$\xi\mapsto\sigma(\xi,\cdot)$\\
$\Jc_\sigma$ & $\sigma$ から決まる一般化複素構造（$\beta$-変換）。$L_\sigma=T^{0,1}\oplus\{\sigma^\sh\xi+\xi:\xi\in T^{*1,0}\}$\\
$A=(T^*)_\sigma$ & 正則 Lie 亜群：アンカー $\sigma^\sh$、括弧 $[a,b]_\sigma=\mathcal L_{\sigma^\sh a}b-\mathcal L_{\sigma^\sh b}a-d\sigma(a,b)$\\
$D=\sum D_i$ & $\sigma$ の退化因子（簡約、単純正規交差または曲面の通常二重点のみ）\\
$T_D=T_X(-\log D)$,\ $B=\Omega^1(\log D)$ & 対数接束、対数余接束。$T_D$ は Lie 亜群、$B=T_D^*$\\
$U$ & $X\setminus D$\\
対数シンプレクティック & $\sigma\in H^0(\wedge^2T_D)$ かつ $\sigma^\sh:B\to T_D$ が同型\\
$L$（直線束として）& 文脈で $\bT_\C$ の固有束と区別する。直線束の方は $E$、$L_0$、$\mathcal O(m)$ 等と書くことも多い（\S\ref{sec:setting}）\\
Poisson 加群 & 平坦 $A$-接続 $\nabla$ をもつ正則ベクトル束 $(E,\nabla)$。一般化正則（GH）束と同じもの（命題 \ref{prop:ident}）\\
$V$ & 局所枠 $e$ に対し $\nabla_\beta e=\beta(V)e$ で定まる正則ベクトル場（Poisson ベクトル場、補題 \ref{lem:flat}）\\
$r_i$ & $D_i$ に沿った対数接続の留数（定数）\\
$C(A;E)$ & $(\wedge^bT\otimes E,\ d_E=[\sigma,\cdot]+V\wedge)$（$A$ の Chevalley--Eilenberg 複体）\\
$\HH^k(E)$ & $H^k(X;C(A;E))$：一般化 Dolbeault コホモロジー（命題 \ref{prop:ident}）\\
$E_1^{a,b}$ & $H^a(X,\wedge^bT\otimes E)$：$\HH^{a+b}$ へ収束するスペクトル系列の $E_1$ 項（$a$ を横断次数と呼ぶ）\\
$\mathrm{DR}_{\log}(E)$ & 対数 de Rham 複体 $(\Omega^\bullet(\log D)\otimes E,\nabla^{\log})$\\
$Q$ & 商複体 $C(A;E)/C(T_D\text{-部分複体})$。$D$ 上に台をもつ\\
$\mu$ & （曲面）$D$ の通常二重点の個数\\
$c,\ \rho,\ b_1$ & $D$ の既約成分数 $c$、$[D_i]\in H^2(X;\C)$ の張る空間の次元 $\rho$、$b_1=c-\rho$\\
$j_1:U\hookrightarrow X',\ j_2:X'\hookrightarrow X$ & 開埋め込み（定理 \ref{thm:main2}）\\
$L_U$ & $U$ 上の $\nabla$-平坦切断の局所系（階数 1、$D_i$ まわりのモノドロミー $e^{2\pi ir_i}$）\\
$P_E(t)$ & $\sum_k\dim\HH^k(E)\,t^k$：Poincaré 多項式\\
$[\,\cdot\,]$（表） & $[\dim\HH^0,\dim\HH^1,\dots,\dim\HH^{2n}]$ の略記\\
\bottomrule
\end{longtable}

%--------------------------------------------------------------
\section{古典の三種の神器：主張・仮定・証明のどこで効くか}\label{sec:classical}
%--------------------------------------------------------------
一般化の目標を明確にするため、古典版の標準的な主張と証明の骨格を整理する。内容は教科書の標準的なもので、新しい結果ではない。

\begin{thm}[Hirzebruch--Riemann--Roch]\tC\label{thm:HRR}
$M$ をコンパクト複素多様体、$E$ を正則ベクトル束とすると
$\chi(M,E)=\sum_q(-1)^q\dim H^q(M,\mathcal O(E))=\int_M\ch(E)\Td(T^{1,0}M).$
\src{Hirzebruch, \emph{Neue topologische Methoden in der algebraischen Geometrie}, Springer 1956（射影代数的な場合）；Atiyah--Singer, Ann.\ Math.\ 87 (1968) 484--530, 531--545, 546--604（一般のコンパクト複素多様体、Dolbeault 複体の指数）。教科書：Huybrechts, \emph{Complex Geometry}, Springer 2005, \S5.1。}
\end{thm}
\begin{hyptab}
コンパクト & コホモロジーが有限次元になる。 & Dolbeault 複体が楕円型なので Fredholm、指数が定義できる。\\\midrule
正則構造 & $\dbar_E^2=0$ が複体を作る条件。 & 複体 $(\Omega^{0,\bullet}(E),\dbar_E)$ の定義。\\\midrule
（結論の形）$\ch(E)$ & 指数が主シンボルの位相的データの積分になる。 & Atiyah--Singer 指数定理。シンボルから $\Td$ が現れ、$E$ は階数ではなく $\ch(E)$ として入る。
\end{hyptab}

\begin{thm}[Serre の消滅定理]\tC\label{thm:Serre}
$X$ を体上の射影的スキーム、$\mathcal O(1)$ を非常に豊富な直線束、$\mathcal F$ を連接層とする。このとき $m_0$ があって、$m\ge m_0$ ならば (a) $H^q(X,\mathcal F(m))=0\ (q>0)$、(b) $\mathcal F(m)$ は大域切断で生成される。
\src{Serre, \emph{Faisceaux alg\'ebriques coh\'erents}, Ann.\ Math.\ 61 (1955) 197--278；現代的証明は Hartshorne, \emph{Algebraic Geometry}, GTM 52, Springer 1977, Theorem III.5.2。}
\end{thm}
\begin{hyptab}
連接性 & 有限生成性により全射 $\bigoplus\mathcal O(-N_i)\twoheadrightarrow\mathcal F$ が取れる。 & $\mathbb P^N$ 上の既知の計算に帰着するため。\\\midrule
豊富性 & ねじり $\mathcal F(m)$ を大きくすると高次コホモロジーが効かなくなる。 & 結論そのもの。\\\midrule
射影的 & $\mathbb P^N$ に閉埋め込みでき、$H^q(\mathbb P^N,\mathcal O(m))$ を Cech 複体で直接計算できる。 & 全射の核 $\mathcal K$ の長完全列に対して $q$ について下降帰納法（$q>\dim X$ で消える）を回すとき。
\end{hyptab}

\begin{thm}[Kodaira--Akizuki--Nakano 消滅]\tC\label{thm:KAN}
$M$ をコンパクト Kähler 多様体、$E$ を正の正則 Hermite 直線束とする。$p+q>n$ ならば $H^q(M,\Omega^p\otimes E)=0$。特に $H^q(M,K_M\otimes E)=0\ (q>0)$。
\src{Kodaira, Proc.\ Nat.\ Acad.\ Sci.\ USA 39 (1953) 1268--1273；Akizuki--Nakano, Proc.\ Japan Acad.\ 30 (1954) 266--272。教科書：Griffiths--Harris, \emph{Principles of Algebraic Geometry}, Wiley 1978, Ch.\ 1 \S2；Huybrechts \S5.2。}
\end{thm}
\begin{hyptab}
コンパクト & (i) Hodge 理論（$H^q\cong\ker\Delta_{\dbar_E}$）、(ii) 部分積分で境界項が出ない。 & 調和表現の存在と、下の積分評価。\\\midrule
Kähler & Kähler の恒等式 $[\Lambda,\dbar]=-i\partial^*$。 & Bochner--Kodaira--Nakano の恒等式 $\Delta_{\dbar_E}=\Delta_{\partial_E}+[i\Theta_E,\Lambda]$ の導出。\\\midrule
正（$i\Theta_E>0$） & 曲率項 $[i\Theta_E,\Lambda]$ が $p+q>n$ で正定値になる。 & $s\in\ker\Delta_{\dbar_E}$ に対し $0=\|\partial_Es\|^2+\|\partial_E^*s\|^2+\langle[i\Theta_E,\Lambda]s,s\rangle$ の最後の項が $>0$ なので $s=0$。
\end{hyptab}

\begin{thm}[Kodaira の埋め込み定理]\tC\label{thm:Kod}
$M$ をコンパクト複素多様体、$L$ を正の正則直線束とすると、$M$ は射影的で、十分大きな $m$ について $|L^m|$ が定める写像 $M\to\mathbb P(H^0(M,L^m)^*)$ は埋め込みである。
\src{Kodaira, Ann.\ Math.\ 60 (1954) 28--48；Griffiths--Harris Ch.\ 1 \S4；Huybrechts \S5.3。}
\end{thm}
\begin{hyptab}
$L$ が正 & 曲率形式 $i\Theta_L$ が Kähler 形式になり、$M$ が Kähler になる。 & Kodaira 消滅を使える状況の確保。\\\midrule
コンパクト & $H^0(M,L^m)$ が有限次元で射影空間への写像が定義できる。 & 写像の構成。\\\midrule
正則切断が十分多い & 標準的証明は点 $x$ での爆発 $\tilde M$、例外因子 $E_x$ を取り、$H^1(\tilde M,\pi^*L^m\otimes\mathcal O(-2E_x))$ の消滅（Kodaira 消滅）から $H^0(L^m)\to L^m\otimes\mathcal O/\mathfrak m_x^2$ の全射性を導く。 & 点の分離と接ベクトルの分離。Riemann--Roch の次元評価を使わずに済むのが要点。
\end{hyptab}

\begin{rem}\tK\ （一般化への示唆）
古典では (1) Hodge 理論（楕円型）、(2) Kähler 恒等式 $\Rightarrow$ BKN 恒等式、(3) 曲率の正値性、(4) 大域切断の多さ、の4点が役割を分担している。一般化複素幾何では、(1) は楕円型複体 $(\wedge L^*\otimes E,d_L)$ として成立する（命題 \ref{prop:ellip}）。(2),(3) の直線束係数版は一般化 Kähler の枠組みではまだ書かれていない。(4) は、一般化正則切断が Poisson 加群の Casimir 切断になる（\S\ref{sec:ghfun}）ため、古典と同じ形では成り立たない。本ノートはこれらを、最も簡単な非自明な例（型が跳ぶ $\beta$-変換）で調べたものである。
\end{rem}

%--------------------------------------------------------------
\section{一般化複素幾何の基礎}\label{sec:gc}
%--------------------------------------------------------------
\subsection{定義}
\begin{defn}[一般化複素構造]\tC\label{def:GC}
$\Jc:\bT\to\bT$ が $\Jc^2=-1$ かつ内積を保つとき一般化概複素構造という。$+i$ 固有束 $L\subset\bT_\C$ が Courant 括弧で閉じるとき $\Jc$ は可積分（一般化複素構造）であるという。
\src{Hitchin, \emph{Generalized Calabi--Yau manifolds}, Q.\ J.\ Math.\ 54 (2003) 281--308, arXiv:math/0209099；Gualtieri, \emph{Generalized complex geometry}, Ann.\ Math.\ 174 (2011) 75--123, arXiv:math/0703298, \S3（型 Def.\ 3.5、標準束 Def.\ 3.7、一般化 Calabi--Yau Def.\ 3.17、Picard 群 Def.\ 3.18。番号は原文で確認）；Gualtieri 博士論文 arXiv:math/0401221。}
\end{defn}
\begin{ex}\tC
複素構造 $I$：$L=T^{0,1}\oplus T^{*1,0}$（型 $n$）。シンプレクティック構造 $\omega$：$L=\{X-i\omega(X)\}$（型 $0$）。
\src{Hitchin 2003（導入）；Gualtieri 2011, \S3（複素構造・シンプレクティック構造の例。具体的な例番号は確認できず。\S4 の Def.\ 4.1 は正則点の定義で、旧稿の「例 4.1--4.2」は存在しない）\tQ。}
\end{ex}
\begin{defn}[型]\tC
$\mathrm{type}(\Jc)_p=\mathrm{codim}_\C\,\pi_T(L_p)$（$\pi_T:\bT\to T$）。型が局所定数な点を正則点という。型は一般に跳ぶ。
\src{Gualtieri 2011, \S3 および Abouzaid--Boyarchenko, \emph{Local structure of generalized complex manifolds}, J.\ Symplectic Geom.\ 4 (2006) 43--62, arXiv:math/0511049。}
\end{defn}
\begin{prop}[楕円型複体]\tC\label{prop:ellip}
$L$ は複素 Lie 亜群（アンカー $\pi_T|_L$）で、$L^*\cong\Lb$ により $\wedge^\bullet L^*$ に Lie 亜群微分 $d_L$ が入る。$(\wedge^\bullet L^*,d_L)$ は楕円型複体で、コンパクト多様体上で一般化 Dolbeault コホモロジー $H^k(\wedge L^*,d_L)$ は有限次元。スピノール束 $K$ を用いると $\wedge^{n-k}\Lb\cdot K\cong U^k$ 上で $d_L$ は $\dbar:U^k\to U^{k-1}$ に対応する。
\src{Gualtieri 2011, Prop.\ 3.13（$(\wedge L^*,d_L)$ は楕円型）、Cor.\ 3.14（コンパクトなら有限次元コホモロジー）。\S\ref{sec:duality} の証明中、$\Jc_\sigma$ の場合には楕円性を主シンボルから直接確かめる（\tP）。}
\end{prop}

\subsection{$\beta$-変換 $\Jc_\sigma$ と一般化正則束}\label{sec:setting}
$X$ を複素多様体、$\sigma$ を正則 Poisson 構造とする。

\begin{prop}\tC\label{prop:Jsigma}
(a) $L_\sigma=T^{0,1}\oplus\{\sigma^\sh\xi+\xi:\xi\in T^{*1,0}\}$ は $\bT_\C$ の極大迷零部分束で、可積分である。対応する一般化複素構造 $\Jc_\sigma$ の型は $p\in X$ で $n-\mathrm{rk}\,\sigma^\sh_p$（$\mathrm{rk}$ は複素階数）。特に $\sigma$ が非退化な点では型 0（シンプレクティック型）、$\sigma$ が零の点では型 $n$（複素型）。
(b) $(\Xi,\xi)\mapsto(\Xi+\sigma^\sh\xi,\xi)$ は複素 Lie 亜群の同型 $T^{0,1}\bowtie(T^{1,0})^*_\sigma\to L_\sigma$ である（左辺は matched pair で、LGSX では $T^{0,1}\bowtie A^{1,0}$ と書く）。
\src{(a) Gualtieri, \emph{Branes on Poisson varieties}, arXiv:0710.2719, Prop.\ 5；Laurent-Gengoux--Stiénon--Xu（以下 LGSX）, \emph{Holomorphic Poisson manifolds and holomorphic Lie algebroids}, IMRN 2008, Art.\ ID rnn088, arXiv:0707.4253, Thm.\ 2.7 と Rem.\ 2.8（そこでは $\Jc_\pi$, $\pi=\pi_R+i\pi_I$ と書く）。型の公式は筆者の整理で、Gualtieri の例（$\mathbb{CP}^2$ 等）と整合する（\tQ）。(b) LGSX Thm.\ 4.21 と Lemma 4.22（$\Jc_{4\pi}$ の因子 4 は $\pi=\pi_R+i\pi_I$ の規格化による。LGSX Prop.\ 3.7）。Gualtieri の論文と LGSX で固有値の符号の取り方が異なる（$\ker(\Jc\mp i)$）が、複素共役で移るので以下の主張には影響しない。}
\end{prop}

\begin{defn}[一般化正則束]\tC
複素ベクトル束 $E$ に平坦 $L$-接続 $\dbar_E:\Gamma(E)\to\Gamma(L^*\otimes E)$（$\dbar_E^2=0$）が与えられたもの。
\src{Gualtieri 2011, \S3（一般化正則束は Thm.\ 3.15 の後に散文で導入され、番号付きの定義はない。複素構造の場合は Example 3.19。Def.\ 3.18 は Picard 群の定義）；Gualtieri, \emph{Branes on Poisson varieties}, \S4.1。}
\end{defn}
\begin{defn}[Poisson 加群]\tC
正則ベクトル束 $E$ と平坦 $A$-接続 $\nabla:\Gamma(E)\to\Gamma(A^*\otimes E)$、すなわち $\nabla_{\xi}(fs)=f\nabla_\xi s+(\sigma^\sh\xi)(f)s$、$[\nabla_a,\nabla_b]=\nabla_{[a,b]_\sigma}$（$a,b$ は正則 1-形式）を満たすもの。ここで $A=(T^*)_\sigma$。
\src{LGSX, Def.\ 4.15 と Lemma 4.16；Gualtieri, \emph{Branes on Poisson varieties}, \S2--4。}
\end{defn}

\begin{prop}[一般化正則束 $=$ Poisson 加群、および Dolbeault モデル]\label{prop:ident}
$(X,\sigma)$ を上の通りとする。
\begin{enumerate}
\item[(a)] $\Jc_\sigma$ に関する一般化正則束は、正則ベクトル束に Poisson 加群構造を与えたものと同じである。\tC\ Gualtieri, \emph{Branes on Poisson varieties}, Prop.\ 8；LGSX Lemma 4.16。直接の証明は下に示す（\tP）。直線束に対しては、Gualtieri 同 Prop.\ 9 により、障害は $H^1(T^{1,0})$ の類 $\sigma\cdot\alpha$ と $H^2_\sigma$ の類で、構造が存在するとき $H^1_\sigma(X)=\HH^1(\mathcal O)$ 上のアフィン空間をなす。\tC
\item[(b)] $L^*=T^{*0,1}\oplus T^{1,0}$ より $\Gamma(\wedge^kL^*\otimes E)=\bigoplus_{a+b=k}A^{0,a}(\wedge^bT\otimes E)$ で、$d_L=\dbar+(-1)^ad_A$、ただし $d_A=[\sigma,\cdot]+V\wedge$ は正則枠の係数に作用する。したがって
\[H^k_{GC}(E)=H^k\bigl(\Gamma(\wedge^\bullet L^*\otimes E),d_L\bigr)=\mathbb H^k\bigl(X;(\wedge^\bullet T\otimes E,\ [\sigma,\cdot]+V\wedge)\bigr)=:\HH^k(E).\]
\tC\ これは LGSX Thm.\ 4.19（$H^*(A,E)=H^*(T^{0,1}\bowtie A^{1,0},E)$、$H^*(A,E)$ は Def.\ 4.17 で層コホモロジーとして定義）。\tP\ 証明の構造は下に示すが、matched pair に関する Mokri の定理（Matched pairs of Lie algebroids, Glasgow Math.\ J.\ 39 (1997) 167--181, doi:10.1017/S0017089500032055。年・DOI は確認、巻頁は記憶\tQ。LGSX Lemma 4.16：複素ベクトル束が $A$-加群であることと $T^{0,1}X\bowtie A^{1,0}$-加群であることは同値。Thm.\ 4.19 の証明は Thm.\ 4.11 と同様とだけ記される）と LGSX Lemma 4.22（再導出していない）に依存する。
\end{enumerate}
\end{prop}
\begin{proof}[(a) の証明（概略）]
($\Leftarrow$) Poisson 加群構造と正則枠 $\{e_s\}$ があるとき、$\nabla_{\bar\Xi}\sum f_se_s=\sum(\bar\Xi f_s)e_s$（正則構造）、$\nabla_\xi\sum f_se_s=\sum((\sigma^\sh\xi)f_s)e_s+f_s\nabla_\xi e_s$（$\xi\in A^{1,0}$）とおく。これは Leibniz 則から強制される唯一の拡張で、$\nabla_\xi e_s$ が正則なので定義可能。平坦性は、$(\bar\Xi,\bar\Upsilon)$ 型は正則構造の可積分性、$(\bar\Xi,\xi)$ 型は「$\nabla_\xi$ は正則切断を正則切断に写す」ことと同値で構成から従い、$(\xi,\eta)$ 型は Poisson 加群の平坦性（正則 $\xi,\eta$ について）を $C^\infty$-線形に拡張したもの（曲率はテンソル）。
($\Rightarrow$) 平坦 $L$-接続を $T^{0,1}$ に制限すると正則構造（Koszul--Malgrange）、混合型の平坦性から $\nabla_\xi$ は正則切断を正則切断に写し、Leibniz 則と平坦性から加群の条件が出る。\qedhere
\end{proof}
\begin{proof}[(b) の証明の構造]
$L\cong T^{0,1}\bowtie A^{1,0}$（命題 \ref{prop:Jsigma}(b)）の Chevalley--Eilenberg 複体は、matched pair の一般論（Mokri）により、$T^{0,1}$ の複体 $C(T^{0,1};\,C(A^{1,0};E))$ の全複体になる。ここで $T^{0,1}$ 方向は $\wedge^bT\otimes E$ の Dolbeault 分解（微分 $\dbar$）で、$A^{1,0}$ 方向は $C(A;E)=(\wedge^bT\otimes E,\ [\sigma,\cdot]+V\wedge)$。2つの微分は Koszul 符号 $(-1)^a$ のもとで反可換。全複体は複体 $(\wedge^\bullet T\otimes E,d_A)$ の細かい分解なので、そのコホモロジーは超コホモロジーである。\qedhere
\end{proof}

\begin{rem}[\tC\ 積分可能性に関する注意]
複素 Lie 亜群 $L=L_\sigma$ 自体を積分する Lie 亜群の概念は、調べた範囲の文献には見つからなかった。LGSX arXiv:0803.2031 の Thm.\ 3.17, 3.22 は、正則 Lie 亜群 $A$ と実 Poisson 構造 $\mathrm{Re}\,\sigma$, $\mathrm{Im}\,\sigma$ の積分可能性の関係を述べるが、$T^{0,1}\bowtie A^{1,0}$ は扱わない（同論文 Rem.\ 4.14 は問題として提起するのみ）。Bailey--Gualtieri, arXiv:1611.03850（題名は\tQ）の Thm.\ 7.1 は、一般化複素構造が弱正則シンプレクティック亜群（Morita 類の概念）に積分されることと、その実 Poisson 構造が積分可能であることの同値を述べる。\tP\ $\mathbb{CP}^2$ の滑らかな3次曲線 $\sigma$ では、$\mathrm{Im}\,\sigma$ の葉は開葉 $U$（等方群 0）と $D$ の各点（等方群 $\mathrm{aff}(1,\C)$、中心が自明）で、Crainic--Fernandes, Ann.\ Math.\ 157 (2003) 575--620（arXiv:math/0105033）Cor.\ 5.2「すべての $N_x$ が自明なら積分可能」により $\mathrm{Im}\,\sigma$ は積分可能。Bailey--Gualtieri の仮定は取得した要約でのみ確認\tQ。
\end{rem}

\subsection{一般化正則切断について}\label{sec:ghfun}
\begin{obs}\tP\label{obs:casimir}
Poisson 加群 $(E,\nabla)$ の一般化正則切断（$\dbar_E s=0$）は、$\HH^0(E)=\{s\in H^0(X,E):\nabla_\xi s=0\ (\forall\xi)\}$、すなわち平坦（Casimir 型）切断である。型 0 の点（$\sigma$ 非退化）の近くでは $\nabla$ は $\sigma^\sh$ により通常の平坦接続に対応し、平坦切断は局所定数になる。
\end{obs}
\begin{proof}
$\HH^0=\ker(d_A:H^0(E)\to H^0(T\otimes E))$ で、$d_A s=(\xi\mapsto\nabla_\xi s)$ である（命題 \ref{prop:ident}(b)）。
\end{proof}
\begin{rem}\tK
これは、型が正則で 0 でない場合（Pal--Poddar の枠組み）にも葉の方向に局所定数となり、一般化正則切断が $M$ から射影空間への写像として埋め込みを与えることはできないことを意味する。Kodaira 埋め込み定理の一般化は、葉空間（あるいは Casimir 切断が分離する商）への写像として定式化し直す必要がある。本ノートでは Kodaira 埋め込みの一般化そのものは扱っていない。
\end{rem}

%--------------------------------------------------------------
\section{対数シンプレクティックな場合：Poisson 加群 $=$ 対数的平坦接続}\label{sec:logsymp}
%--------------------------------------------------------------
\subsection{設定}
$X$ をコンパクト複素多様体（$\dim_\C X=n=2k$）、$\sigma$ を正則 Poisson 構造、$D\subset X$ を簡約な因子（単純正規交差、曲面では通常二重点のみも許す）とする。
$T_D=T_X(-\log D)$、$B=\Omega^1(\log D)$。

\begin{defn}[対数シンプレクティック]\tC
$\sigma\in H^0(\wedge^2T_D)$ かつ $\sigma^\sh:B\to T_D$ が同型であるとき、$\sigma$ は $D$ に沿って対数シンプレクティックであるという（同値：$\omega:=\sigma^{-1}\in H^0(\Omega^2(\log D))$ が非退化）。このとき $\mathrm{div}(\sigma^k)=D$、すなわち $-K_X=\mathcal O(D)$。
\src{Goto, \emph{Unobstructed deformations of generalized complex structures induced by $C^\infty$ logarithmic symplectic structures}, arXiv:1501.03398（掲載誌・巻頁は未確認\tQ）；Gualtieri--Li, \emph{Symplectic groupoids of log symplectic manifolds}, IMRN 2014, arXiv:1206.3674；Pym, \emph{Constructions and classifications of projective Poisson varieties}, Lett.\ Math.\ Phys.\ 108 (2018) 573--632（確認）。Gualtieri--Li の掲載誌と Goto の掲載情報は未確認\tQ。}
曲面の場合これは、$\sigma\in H^0(-K_X)$ で零因子が簡約かつ通常二重点しか持たない、ということである。局所的に $\sigma=h\,\partial_1\wedge\partial_2$、$h=zu$（滑らかな点）または $h=z_1z_2u$（二重点）、$u$ は単元。
\end{defn}

\subsection{主張}
\begin{thm}[Poisson 加群 $=$ 対数的平坦接続]\tP\label{thm:main1}
$(X,D,\sigma)$ を対数シンプレクティック（任意の偶数次元、$D$ は単純正規交差）、$E$ を Poisson 加群を持つ正則直線束とする。
\begin{enumerate}
\item[(a)] 局所枠 $e$ に対し $\nabla_\beta e=\beta(V)e$ で正則ベクトル場 $V$ が定まり、$[\sigma,V]=0$（Poisson ベクトル場）。枠の取り替え $e\mapsto ge$ で $V$ は $X_{\log g}$ 型のハミルトン項だけ変わる。
\item[(b)] $V$ は自動的に対数的：$V\in T_D$。したがって $A$-加群構造は一意に $B$-加群構造に延び、同型 $\sigma^\sh:B\cong T_D$ を通して、対数極をもつ平坦接続 $\nabla^{\log}:E\to\Omega^1(\log D)\otimes E$ と同じものになる。
\item[(c)] $\nabla^{\log}$ の $D_i$ に沿った留数は定数 $r_i\in\C$ で、$H^2(X;\C)$ において $c_1(E)=-\sum_ir_i[D_i]$。
\item[(d)] $D$ の点の近くで、$\nabla^{\log}e=\sum_ir_i\,\dlog z_i\,e$ となる枠が取れる。
\item[(e)] 部分複体 $C(T_D;E)=(\wedge^bT_D\otimes E,d_E)\subset C(A;E)$ は、$\wedge^b\sigma^\sh$ により対数 de Rham 複体 $\mathrm{DR}_{\log}(E)$ と同型。
\end{enumerate}
\end{thm}
\begin{hyptab}
対数シンプレクティック（$\sigma^\sh:B\cong T_D$） & $B$ 上の Lie 亜群構造 $(T_D^*)_\sigma$ と $T_D$ を同一視し、Poisson 加群を通常の平坦接続（対数極つき）として扱うため。 & 補題 \ref{lem:iso}、定理 (b),(e)。\\\midrule
$D$ が簡約・単純正規交差 & 局所方程式 $f=z_1\cdots z_s$ の各因子から極の位数を独立に評価するため。 & 補題 \ref{lem:log}（$V(f)\in(f)$ から $V(z_i)\in(z_i)$ を導く）。\\\midrule
$X$ コンパクト & $c_1(E)$ と留数の関係（大域的なコホモロジー類の等式）のため。 & (c)。留数が定数であることは局所的。\\\midrule
$\sigma$ が Poisson（$[\sigma,\sigma]=0$）& $A$ が Lie 亜群になり、$d_E^2=0$ となる。 & 補題 \ref{lem:flat}。
\end{hyptab}

\begin{lem}[平坦性 $\Leftrightarrow$ Poisson ベクトル場]\tP\label{lem:flat}
局所枠 $e$ について $\Theta(\beta):=\nabla_\beta e/e$ とおくと、$\nabla_a\nabla_be=(\sigma^\sh a)(\Theta b)e+\Theta(a)\Theta(b)e$。よって平坦性 $[\nabla_a,\nabla_b]=\nabla_{[a,b]_\sigma}$ は
$(\sigma^\sh a)(\Theta b)-(\sigma^\sh b)(\Theta a)=\Theta([a,b]_\sigma)$、すなわち $\Theta$ が Lichnerowicz 複体の 1-コサイクル：$\Theta(\beta)=\beta(V)$、$[\sigma,V]=0$ と同値。
局所座標で $\sigma=h\partial_1\wedge\partial_2$、$a=\Theta(dz_1)$、$b=\Theta(dz_2)$ のとき条件は $h(a_1+b_2)=h_1a+h_2b$ すなわち $V(h)=h\,\mathrm{div}V$。微分 $d_E$ は $C^0$ で $d_E(fe)=(-\sigma^\sh(df)+fV)e$、$c=X\otimes e$、$X=p\partial_1+q\partial_2$ に対し
\[(d_Ec)(dz_1,dz_2)=h(p_1+q_2)-h_1p-h_2q+qa-pb.\]
\end{lem}
\begin{lem}[$V$ は対数的]\tP\label{lem:log}
$D$ の局所方程式を $f=z_1\cdots z_s$ とする。$L_V\sigma=0$ から $L_V\sigma^k=0$。$\sigma^k=f'\partial_1\wedge\cdots\wedge\partial_n$（$f'=f\cdot$単元）なので $V(f')-f'\,\mathrm{div}V=0$、よって $V(f)\in(f)$。$V(f)/f=\sum_iV(z_i)/z_i$ で、第 $i$ 項の極は $\{z_i=0\}$ のみ。全体が正則なので各 $V(z_i)/z_i$ が正則、すなわち $V(z_i)\in(z_i)$。
\end{lem}
\begin{lem}[$\sigma^\sh:B\to T_D$ は Lie 亜群の同型]\tP\ （標準）\label{lem:iso}
$\sigma\in\wedge^2T_D$ は Lie 亜群 $T_D$ 上の Poisson 構造で非退化。$\omega=\sigma^{-1}$ が $d_{T_D}$-閉であることは、$U$ 上で $\omega$ が閉（Poisson 構造の逆）で、$\Omega^3(\log D)$ の切断が稠密開集合上で 0 なら 0 であることから従う。閉 2-形式による非退化 Poisson 構造では $\sigma^\sh$ が括弧を保つ。ゆえに $(T_D^*)_\sigma=B\cong T_D$。
\end{lem}
\begin{lem}[留数・局所標準形・$c_1$]\tP\ /\ \tC\label{lem:res}
$\nabla^{\log}e=\theta e$、$\theta\in\Omega^1(\log D)$ は平坦性から閉。$D_i$ に沿った留数 $\theta_i|_{D_i}$ は、$d\theta$ を $D_i$ に制限すると定数であることが分かる。$\theta-\sum r_i\dlog z_i$ は閉正則 1-形式なので $=dg$、枠 $e'=e^{-g}e$ で $\nabla^{\log}e'=\sum r_i\dlog z_i\,e'$（(d)）。$c_1(E)=-\sum r_i[D_i]$ は Esnault--Viehweg の留数公式（\tC\ Esnault--Viehweg, \emph{Logarithmic de Rham complexes and vanishing theorems}, Invent.\ Math.\ 86 (1986) 161--194, \S2；Deligne, \emph{\'Equations diff\'erentielles \`a points singuliers r\'eguliers}, LNM 163, Springer 1970, Ch.\ II。符号の確認：標準切断が平坦な $\mathcal O(D)$ の留数は $-1$、$c_1=[D]$）。
同じ $c_1$ をもつ加群構造の全体は $H^0(\Omega^1(\log D))=\ker(\C^c\to H^2(X;\C))$ 上のトーサー（Deligne）で、次元は $c-\rho$。$D$ が既約なら一意、可約なら留数ベクトルのアフィン族になる。
\end{lem}
\begin{proof}[定理 \ref{thm:main1} の証明]
(a) 補題 \ref{lem:flat}。(b) 補題 \ref{lem:log} で $\Theta$ が $B$ に延び、平坦性は稠密開集合上の等式なので延長後も成り立つ。補題 \ref{lem:iso} と合わせる。(c),(d) 補題 \ref{lem:res}。(e) 補題 \ref{lem:iso}。
\end{proof}

%--------------------------------------------------------------
\section{最も簡単な非自明な例：$\mathbb{CP}^2$ と滑らかな3次曲線}\label{sec:cp2}
%--------------------------------------------------------------
$F$ を $\mathbb C^3$ 上の滑らかな3次斉次多項式、$\Pi^{ij}=\epsilon_{ijk}\partial_kF$ とする。$\Pi$ は Euler ベクトル場と縮約すると 0 になる斉次 Poisson 構造（ウェイト 0）で、$\mathbb{CP}^2$ 上の正則 Poisson 構造 $\sigma\in H^0(\wedge^2T)=H^0(\mathcal O(3))$ を定める。$\sigma$ は $D=\{F=0\}$（滑らかな3次曲線）で 1 位に消える。標準的な構成である（出典の同定は\tQ；$D$ での消え方は局所計算で確認した）。以下、$E=\mathcal O(m)$ とする。

\subsection{どの $\mathcal O(m)$ が一般化正則か}
\begin{prop}\tP\label{prop:cp2mod}
すべての $m\in\Z$ に対し $\mathcal O(m)$ はただ一つの Poisson 加群構造をもち、その対数接続の留数は $r=-m/3$ である。
\end{prop}
\begin{proof}
定理 \ref{thm:main1}(c) より $r=-m/3$（$[D]=3H$）。存在：$\mathcal O(m)$ の Atiyah 類が $H^1(\Omega^1(\log D))$ で消えることと $c_1\in\C[D]$ が同値（完全列 $H^0(\mathcal O_D)\to H^1(\Omega^1)\to H^1(\Omega^1(\log D))$）で、$m[H]\in\C\cdot3[H]$。曲率は $H^0(\Omega^2(\log D))=\C\omega$ の元で留数 $dr=0$ だから平坦。一意性：2つの対数接続の差は $H^0(\Omega^1(\log D))=0$（連結写像 $1\mapsto[D]=3H\ne0$）。\qedhere
\end{proof}
別の見方：$\Pi$ の Casimir $F$ に対し、$F^{m/3}$ は $U=\mathbb{CP}^2\setminus D$ 上で $\mathcal O(m)$ の（多価の）平坦切断を与える。$\mathcal O(m)|_U$ は階数 1 の局所系 $L_m$ で、$D$ まわりのモノドロミーは $e^{-2\pi im/3}$。$\pi_1(U)=\Z/3$ に対応して周期 3 が現れる。\tP\ 局所的に $\nabla e=-\tfrac m3\dlog z\,e$。

\subsection{結果}
\begin{thm}\tP\label{thm:cp2}
$\mathcal O(m)$ の一般化 Dolbeault コホモロジー
\[\HH^*(\mathcal O(m))=\mathbb H^*\bigl(\mathbb{CP}^2;(\wedge^\bullet T(m),d_\sigma)\bigr)\]
の次元 $[\dim\HH^0,\dots,\dim\HH^4]$ は
\begin{center}\small
\begin{tabular}{lll}\toprule
$m$ & 次元 & 解釈\\\midrule
$m\in3\Z_{\ge0}$ & $[1,0,2,0,0]$ & $H^*(U;\C)$。$\HH^0=\C F^{m/3}$ は Casimir 切断\\
$m\in3\Z_{\le-6}$ & $[0,0,2,0,1]$ & $H^*_c(U;\C)$\\
上記以外（$m=-3$ を含む）& $[0,0,3,0,0]$ & $H^*(U;L_m)$（$m\notin3\Z$）、$m=-3$ は特別\\\bottomrule
\end{tabular}
\end{center}
全 Euler 標数は常に $3=\chi(\mathbb{CP}^2)$。
\end{thm}
\begin{hyptab}
$D$ が滑らか・$\sigma$ が $D$ に沿って 1 位に消える & 商複体 $Q$ の微分が同型 $d\sigma|_D:N_D\to\wedge^2T|_D$ の定数倍になる。 & 証明の手順 (3)。$m\ne-3$ で $Q$ が非輪状。\\\midrule
$\pi_1(U)=\Z/3$（Milnor ファイバー $\simeq$ $S^2$ の8個の束） & $U$ のコホモロジーを具体的に決める。 & 手順 (2) の表の右辺。\\\midrule
$m\ne-3$ & $Q$ の微分 $-(1+m/3)\cdot(d\sigma|_D)$ が零写像でない。 & $m=-3$ だけ別扱い（手順 (4)）。
\end{hyptab}
\begin{proof}[証明の手順]
(1) 命題 \ref{prop:cp2mod} と定理 \ref{thm:main1}(e) により、部分複体 $C_{\log}=(\wedge^bT_D\otimes\mathcal O(m))\cong\mathrm{DR}_{\log}$。
(2) 局所枠で $\nabla^{\log}e=r\,\dlog z\,e$（$r=-m/3$）。$D$ の点での茎では、対数 de Rham 複体は単項式 $z^ne$ に $(n+r)=(n-m/3)$ 倍で作用する。茎のコホモロジーは $m\in3\Z_{\ge0}$ で $[1,1]$、それ以外で 0。ゆえに $m\in3\Z_{\ge0}$ では自然な写像 $\mathrm{DR}_{\log}\to Rj_*L_m$ が茎ごとに擬同型（$\mathbb H(C_{\log})=H^*(U;\C)$）、それ以外では $j_!L_m\to\mathrm{DR}_{\log}$ が擬同型（$\mathbb H(C_{\log})=H_c^*(U;L_m)$）。$U$ の位相：$F=1$ の Milnor ファイバーは $\mu=8$ 個の $S^2$ の束で $U=\{F=1\}/\mu_3$、$H^2=\C^8=2+3+3$（自明指標・非自明指標2つ）。$U$ はアファインなので $H^k=0\ (k>2)$。よって $H^*(U)=[1,0,2,0,0]$、非自明な指標の $H^*(U;L)=[0,0,3,0,0]$、Verdier 双対性で $H^*_c(U;\C)=[0,0,2,0,1]$。
(3) 商 $Q=C(A;\mathcal O(m))/C_{\log}$ は $D$ 上の複体で、$Q^1=N_D(m)=\mathcal O_D(m+3)$、$Q^2=\wedge^2T(m)|_D=\mathcal O_D(m+3)$。補題 \ref{lem:flat} の公式で $z=0$ とすると $d_Q=-u(1-r)p$、不変な形では $d_Q=-(1+m/3)\cdot(d\sigma|_D)$。$d\sigma|_D$ は同型なので、$m\ne-3$ のとき $Q$ は非輪状で $\HH^*(\mathcal O(m))=\mathbb H^*(C_{\log})$。
(4) $m=-3$：$E_1^{a,b}=H^a(\wedge^bT(-3))$ は $(2,0),(1,1),(0,2)$ にそれぞれ 1 次元（Bott）、すべて全次数 2 なので以後の微分は零で $[0,0,3,0,0]$。符号の確認は不要。
\end{proof}
\begin{rem}[\tV]
上の表は、$m=-15,\dots,15$ について、Fermat 型と Hesse 型の3次曲線を使い、スペクトル系列 $E_1^{a,b}=H^a(\wedge^bT(m))$ の微分 $d_1$ を（$m\ge-2$ では大域切断の複体として、$m\le-4$ では Laurent 単項式による Cech 計算として）素数 $32003$ を法とする厳密な線形代数で直接計算した結果と一致する（スクリプト：\texttt{cp2/poisson.py, table.py, neg\_check.py, predict.py}）。
\end{rem}
\begin{rem}[\tC\ 関連]
Hitchin, \emph{Deformations of holomorphic Poisson manifolds}, arXiv:1105.4775（Moscow Math.\ J.\ 12 (2012) 567--591）の (11) 式 $0\to T^*\xrightarrow{\sigma}T\to T|_D\to0$ と Prop.\ 8 が関連するが、$\mathbb{CP}^2$ のこの結果そのものは同論文にも Gualtieri の論文にも書かれていない（調べた範囲）。
\end{rem}

\begin{obs}[「$m$ への依存」の見え方]\tP
全 Euler 標数は $m$ によらない（位相不変量）。$m$ への依存は Poincaré 多項式、特に周期 3 の跳びと $\HH^0$（Casimir 切断 $F^{m/3}$）にのみ現れる。$m\gg0$ でも $\HH^2$ は 2 または 3 次元のまま残るので、全次数での Serre 消滅は成り立たない。
\end{obs}

%--------------------------------------------------------------
\section{主定理（曲面）：対数シンプレクティック曲面の一般化 Dolbeault コホモロジー}\label{sec:main2}
%--------------------------------------------------------------
\subsection{主張}
$X$ をコンパクト複素曲面、$\sigma\in H^0(-K_X)$ は簡約な零因子 $D=\bigcup D_i$（既約成分 $D_i$、$D$ の特異点は通常二重点のみ、その個数 $\mu$）をもち、$E$ は Poisson 加群をもつ直線束とする。留数 $r_i$（各分枝に沿って定数）は定理 \ref{thm:main1}(c) による。（二重点の枝は別の枝と見なす。）
$I_*=\{i:r_i\in\Z_{\le0}\}$、$I_!=\{i:r_i\notin\Z_{\le0}\}$、$X'=X\setminus\bigcup_{i\in I_!}D_i$、$j_1:U\hookrightarrow X'$、$j_2:X'\hookrightarrow X$ とし、
\[\mathcal F:=j_{2!}\,Rj_{1*}L_U\]
（$I_*$ の成分では $Rj_*$、$I_!$ の成分では零拡張）とおく。

\begin{thm}[曲面の主定理]\tP\label{thm:main2}
上の設定で、すべての $r_i\neq1$ と仮定する。このとき
\[\HH^k(E)\;\cong\;H^k(X;\mathcal F)\ \oplus\ \C^{\mu}\ (k=2\text{ のみ}).\]
特に
\begin{enumerate}
\item[(i)] すべての $r_i\notin\Z_{>0}$ なら $\HH^*(E)=H^*(U;L_U)\oplus\C^\mu[-2]$、
\item[(ii)] すべての $r_i\notin\Z_{\le0}$ なら $\HH^*(E)=H^*_c(U;L_U)\oplus\C^\mu[-2]$、
\item[(iii)] 全 Euler 標数 $\sum(-1)^k\dim\HH^k(E)=\chi(U)+\mu=\chi(X)$（$E$ によらない）。
\end{enumerate}
\end{thm}
\begin{hyptab}
曲面、$\sigma\in H^0(-K_X)$、$D$ 簡約・通常二重点のみ & 局所標準形 $h=zu$ または $z_1z_2u$ を使い、商複体 $Q$ を明示的に計算するため。 & 補題 \ref{lem:Q}、命題 \ref{prop:moser}、補題 \ref{lem:conn}。\\\midrule
$r_i\ne1$（全分枝）& $r_i=1$ だと $Q$ の微分 $-u(1-r)p$ が零になり、$Q$ のコホモロジー層が無限次元になって簡約が破れる。 & 補題 \ref{lem:Q}。（$r_i=1$ は $E=K_X$ のときの特別扱い、命題 \ref{prop:K}。）\\\midrule
留数が正整数でない／非正整数でない & 対数複体と $Rj_*L_U$、$j_!L_U$ の比較（Deligne）が擬同型になる条件。 & 補題 \ref{lem:stalk}。\\\midrule
$X$ コンパクト & $R\Gamma(X;j_!L_U)=R\Gamma_c(U;L_U)$ と Euler 標数の計算。 & (ii),(iii)。\\\midrule
$-K_X=\mathcal O(D)$（反標準曲線） & $\chi(D)=\mu$ を導く（$\omega_D=\mathcal O_D$、$p_a(D)=1$）。 & (iii)。
\end{hyptab}

\subsection{証明}
\begin{lem}[対数複体の茎と Deligne の比較]\tP\ /\ \tC\label{lem:stalk}
局所枠（定理 \ref{thm:main1}(d)）で、$D_I$（$|I|=s$）の点での $\mathrm{DR}_{\log}$ の茎は、可換作用素 $z_i\partial_i+r_i\ (i\in I)$ の Koszul 複体と、残りの変数の通常の de Rham 複体のテンソル積である。単項式 $z^a$ 上では $z_i\partial_i+r_i=a_i+r_i$（スカラー）。したがって茎のコホモロジーは、すべての $i\in I$ で $-r_i\in\mathbb N$（すなわち $r_i\in\Z_{\le0}$）のときに限り $\wedge^*\C^s$（次元 $2^s$、Euler 標数 0）で、そうでなければ 0。一方 $Rj_*L_U$ の同じ点での茎は、すべての $e^{2\pi ir_i}=1$（$r_i\in\Z$）のときに限り $\wedge^*\C^s$。ゆえに、
\begin{itemize}
\item 自然な写像 $\mathrm{DR}_{\log}\to Rj_*L_U$ は、正整数の $r_i$ が一つもないとき擬同型（Deligne, LNM 163, Ch.\ II \S3、番号は\tQ。ここでは直接計算で確認）。
\item 零拡張 $j_!L_U\to\mathrm{DR}_{\log}$（平坦切断の零拡張は正則）は、$r_i\in\Z_{\le0}$ となる $i$ が一つもないとき擬同型。
\item 混合の場合、$X'$ の上では $\mathrm{DR}_{\log}\simeq Rj_{1*}L_U$、$\bigcup_{I_!}D_i$ の上では茎が 0 なので、$\mathrm{DR}_{\log}\simeq\mathcal F$。
\end{itemize}
係数 $1/(a_i+r_i)$ による割り算は、有限個の $a$ を除き $|a_i+r_i|\ge c>0$ なので収束を壊さない。
\end{lem}
\begin{lem}[商複体 $Q$]\tP\ \tV\label{lem:Q}
$Q:=C(A;E)/C(T_D;E)$ は $D$ 上に台をもち、次数 1, 2 のみ。$Q^1=(T/T_D)\otimes E$、$Q^2=(\wedge^2T/\wedge^2T_D)\otimes E$。
\begin{itemize}
\item 滑らかな点（$h=zu$）：$Q^1=\partial_z\otimes\mathcal O_D$、$Q^2=\partial_z\wedge\partial_w\otimes\mathcal O_D$、$d_Q(p)=-u(1-r)\,p|_D$。よって $r\ne1$ のとき非輪状。
\item 二重点（$h=z_1z_2u$）：$Q^1=\mathcal O_{D_1}\partial_1\oplus\mathcal O_{D_2}\partial_2$、$Q^2=\mathcal O_D$、$d_Q(p,q)=-u[(1-r_1)z_2p+(1-r_2)z_1q]\bmod z_1z_2$。$r_1,r_2\ne1$ なら単射で像は極大イデアル $\mathfrak m_D$ なので $H^1(Q)_p=0$、$H^2(Q)_p=\mathcal O_D/\mathfrak m_D=\C$。
\end{itemize}
ゆえに $r_i\ne1$ のとき、導来圏で $Q\simeq\bigoplus_{\text{二重点 }p}\C_p[-2]$。
導出：補題 \ref{lem:flat} の公式を $z=0$ に制限する。$a|_D=\theta_2h$、$b|_D=-\theta_1u$、$\theta_1|_D=r$（定理 \ref{thm:main1}(d)）を代入。（\tV\ ランダムな多項式 $u,g,p,q$ と留数で、滑らかな点・二重点とも各20回、$d_1d_0=0$ と $d_Q$ の公式を厳密算術で確認。スクリプト \texttt{cp2/ncd\_cochain\_check.py}。）
\end{lem}
\begin{prop}[二重点での対数 Moser の標準形]\tP\label{prop:moser}
二重点 $p$ の近傍で $\omega=\sigma^{-1}=u^{-1}\dlog z_1\wedge\dlog z_2$、$c:=u(0)$ とする。$D$ の各枝を保ち $p$ を固定する双正則写像 $\phi$ で $\phi^*\omega=c^{-1}\dlog z_1\wedge\dlog z_2$ となるものが存在する。$c$ は局所不変量で、一般には 1 にできない。
\end{prop}
\begin{proof}
$\omega_0=\dlog z_1\wedge\dlog z_2$、$k=1/u-1/c$（$k(0)=0$）、$\omega_t=(c^{-1}+tk)\omega_0$ とおく。$\omega_t$ は閉で、各 $t\in[0,1]$ で $p$ の近くで非退化。$k=\sum k_{ab}z_1^az_2^b$ に対し
\[\alpha=\sum_{a\ge1,b\ge0}\tfrac{k_{ab}}{a}z_1^az_2^b\,\dlog z_2-\sum_{b\ge1}\tfrac{k_{0b}}{b}z_2^b\,\dlog z_1\]
は収束する対数 1-形式で $d\alpha=k\omega_0$（$d(z^{ab}\dlog z_2)=az^{ab}\omega_0$、$d(z_2^b\dlog z_1)=-bz_2^b\omega_0$）。$X\mapsto\iota_X\omega_t$ は $T_D\to\Omega^1(\log D)$ の同型なので、$\iota_{X_t}\omega_t=-\alpha$ を満たす $X_t\in T_D$ が一意に定まり、各枝に接し $p$ で 0 になる。その流れ $\phi_t$ は $D$ を保ち、$\tfrac{d}{dt}\phi_t^*\omega_t=\phi_t^*(d\iota_{X_t}\omega_t+k\omega_0)=\phi_t^*(-d\alpha+k\omega_0)=0$。ゆえに $\phi_1^*\omega_1=\omega_{c^{-1}}$ 型。\qedhere
\end{proof}
\begin{lem}[連結射の消滅]\tP\ \tV\label{lem:conn}
完全三角 $C(T_D;E)\to C(A;E)\to Q\to C(T_D;E)[1]$ の連結射の $p$ 成分 $e_p\in H^3_p(X;\mathrm{DR}_{\log})$ は、$(r_1,r_2)\ne(1,1)$ のとき 0 である。
\end{lem}
\begin{proof}[概略]
命題 \ref{prop:moser} と補題 \ref{lem:res} により $\sigma=c\,z_1z_2\partial_1\wedge\partial_2$、$\nabla^{\log}e=(r_1\dlog z_1+r_2\dlog z_2)e$ としてよい。すべての複体はトーラスウェイト $\gamma\in\Z^2$ で次数付けられ、微分はウェイトを保つ（$\lambda_i=\gamma_i+r_i$）。局所コホモロジーは Laurent 尾 $H^2_p(\mathcal O)=\mathrm{span}\{z^{-a}:a\ge1\}$ で計算できる。(a) $H^3_p(\mathrm{DR}_{\log})$ は、$\lambda_1=\lambda_2=0$（$\gamma=-r$）のウェイトにのみ現れ、$r_1,r_2$ が正整数のときに限り非零。(b) $H^2_p(Q)=\C$ は単一ウェイト $\gamma=(-1,-1)$ にあり（$\mathcal O_D$ の局所コホモロジーに余分な類 $H^1_p(\mathcal O_D)\ni[c_1\ne c_2]$ が必要）、$(-1,-1)=-r$ は $r_1=r_2=1$ のときに限る。連結射はウェイトを保つので、$(r_1,r_2)\ne(1,1)$ のとき $e_p=0$。ウェイトごとの Koszul 複体の計算と合わせ、収束の問題はない（すべての写像が単項式上で対角）。
\tV\ スクリプト \texttt{cp2/ncd\_local\_weights.py}：ウェイト $|\gamma_i|\le8$、$r_i\in\{-3..3,\pm\frac12,\frac13,-\frac23,\frac52\}$ の 78 組（すべて $r_i\ne1$）で、茎のコホモロジー $[0,0,1]$、局所コホモロジーの錐 $=H^2_p$（1 クラス）、$H(\mathrm{LOG})\to H(\mathrm{FULL})$ の単射性（連結射が 0）を厳密算術で確認。違反 0。$r=(1,1)$ では連結射は非零。
\end{proof}
\begin{proof}[定理 \ref{thm:main2} の証明]
定理 \ref{thm:main1}、補題 \ref{lem:stalk} により $C(T_D;E)\simeq\mathrm{DR}_{\log}\simeq\mathcal F$。補題 \ref{lem:Q}、\ref{lem:conn} により $C(A;E)\simeq\mathrm{DR}_{\log}\oplus\bigoplus_p\C_p[-2]$（$D^b(\C_X)$ 内）。$R\Gamma(X;\cdot)$ を取る。
(i),(ii)：$\mathcal F=Rj_*L_U$ または $j_!L_U$、$X$ コンパクトより $R\Gamma(X;j_!L_U)=R\Gamma_c(U;L_U)$。
(iii)：$\chi(X;\mathcal F)=\sum_S\chi_c(S)\chi(\mathcal F_x)$ で、$D$ 上の茎の Euler 標数は 0（$\wedge^*\C^s$ または 0）なので $=\chi_c(U)=\chi(U)=\chi(X)-\chi(D)$。$D$ は反標準なので $\omega_D=\mathcal O_D$、$\chi(\mathcal O_D)=0$ から $\sum_j(2-2\tilde g_j)=2\mu$（$\tilde D_j$ は正規化）、よって $\chi_{\rm top}(D)=\mu$、$\chi(U)+\mu=\chi(X)$。（単純正規交差の場合の別証：$\chi(U)=\int c_2(T_X(-\log D))$。）
\end{proof}
\begin{prop}[$r_i=1$：$E=K_X$]\tP\label{prop:K}
$E=K_X$ なら $\wedge^bT\otimes K_X=\Omega^{2-b}$ で、$E_1^{a,b}=H^{2-b,a}(X)$（全次数 $2+a-p$、$p=2-b$）。$h^{p,q}(X)=0\ (p\ne q)$ なら $E_1$ は全次数 2 のみに集中し、微分は全次数を 1 上げるので退化し、$\HH^k=0\ (k\ne2)$、$\dim\HH^2=\sum_ph^{p,p}=\chi(X)$（$\mathbb{CP}^2$ で 3、$\mathbb{CP}^1\times\mathbb{CP}^1$ で 4）。
\end{prop}

\subsection{$\mathbb{CP}^1\times\mathbb{CP}^1$ での確認}
$\sigma=\Pi_F$、$F\in H^0(\mathcal O(2,2))$（$-K=\mathcal O(2,2)$）。\tP\ 滑らかな $D$ のとき、$\mathcal O(a,b)$ が Poisson 加群をもつのは $a=b=m$ のときに限り、留数は $-m/2$、構造は一意（\S\ref{sec:cp2} と同じ議論）。\tV\ $m=-9,\dots,8$ の厳密計算（スクリプト \texttt{cp2/p1p1.py}）：
\begin{center}\small
\begin{tabular}{lll}\toprule
$D$ & $m$ 偶数 $\ge0$ / 偶数 $\le-4$ & $m$ 奇数、$m=-2$\\\midrule
滑らかな $(2,2)$ 曲線 & $[1,0,3,0,0]$ / $[0,0,3,0,1]$ & $[0,0,4,0,0]$\\
横断する2つの $(1,1)$ 曲線（$\mu=2$）& $[1,1,4,0,0]$ / $[0,0,4,1,1]$ & $[0,0,4,0,0]$\\
4本の線 $x_0x_1y_0y_1$（$\mu=4$）& $[1,2,5,0,0]$ / $[0,0,5,2,1]$ & $[0,0,4,0,0]$\\
$(1,1)+(1,0)+(0,1)$（$\mu=3$）& $[1,1,4,0,0]$ / $[0,0,4,1,1]$ & $[0,0,4,0,0]$\\\bottomrule
\end{tabular}
\end{center}
周期 2、Serre 双対性 $m\leftrightarrow-m-4$、$\chi=4$ は常に成立。可約な $D$ では Poisson 加群構造は一意でなく（$H^0(\Omega^1(\log D))\ne0$）、上の表は成分ごとに留数が等しい対称な構造について計算している（補題 \ref{lem:res}）。

\subsection{定理 \ref{thm:main2} の検証}
\tV\ 定理 \ref{thm:main2} とその閉公式（\S\ref{sec:RR} の定理 \ref{thm:B}）を、次の8つの例の合計 196 個の $m$ について、上の厳密計算の結果と比較した（スクリプト \texttt{cp2/ncd\_global\_check.py}）：$\mathbb{CP}^2$ 上の滑らかな3次曲線、結節3次曲線（$\mu=1$）、円錐＋割線（$\mu=2$）、三角形（$\mu=3$）（$m=-15..15$）、および上の $\mathbb{CP}^1\times\mathbb{CP}^1$ の4例（$m=-9..8$）。食い違い 0。$\mathbb{CP}^2$ の尖点3次、円錐＋接線（tacnode）、3線共点、非簡約 $D$ はこの定理の対象外で、\S\ref{sec:local} で扱う。

%--------------------------------------------------------------
\section{Serre 型の定理}\label{sec:serre}
%--------------------------------------------------------------
\subsection{古典との対応}
古典の Serre の消滅定理（定理 \ref{thm:Serre}）の各要素に対応する一般化版は次の通りである（\tK\ 対応づけは筆者の整理）。
\begin{center}\small
\begin{tabular}{p{4.6cm}p{9.8cm}}\toprule
古典 & 本ノートでの対応物\\\midrule
連接層 $\mathcal F$ & Lie 亜群複体の各項 $K^b=\wedge^bT\otimes E$（$E$ は一般化正則束）。連接である。\\
豊富な $\mathcal O(1)$ & 豊富な正則直線束 $L_0$ で、かつ Poisson 加群（一般化正則）でもあるもの。Poisson 構造は $C(A;\cdot)$ が複体になるためだけに必要で、正値性には関与しない。\\
射影的 & $X$ 射影的（$L_0$ が存在するため）。\\
$H^q(\mathcal F(m))=0\ (q>0)$ & $E_1^{a,b}=H^a(X,\wedge^bT\otimes E\otimes L_0^m)=0$ （\emph{$a>0$：横断（Dolbeault）次数}）。\\
結論：$H^*(\mathcal F(m))$ が次数 0 に集中 & $\HH^*$ は大域切断の有限次元複体 $(\Gamma(\wedge^\bullet T\otimes E\otimes L_0^m),d)$ のコホモロジー。\\\bottomrule
\end{tabular}
\end{center}

\subsection{任意次元の Serre 型定理}
\begin{thm}[Serre 型定理（$\beta$-変換）]\tP\label{thm:A1}
$X$ をコンパクト複素射影多様体（$\dim_\C X=n$）、$\sigma$ を任意の正則 Poisson 構造（$[\sigma,\sigma]=0$、非退化性や対数条件は不要）、$E$ を一般化正則ベクトル束（Poisson 加群）、$L_0$ を Poisson 加群構造をもつ豊富な正則直線束とする。このとき、$X,\sigma,E,L_0$ のみに依存する $m_0$ が存在して、$m\ge m_0$ ならば
\begin{enumerate}
\item[(a)] すべての $a>0$、すべての $b$ で $H^a(X,\wedge^bT\otimes E\otimes L_0^m)=0$、
\item[(b)] $\HH^k(E\otimes L_0^m)\cong H^k\bigl(\Gamma(X,\wedge^\bullet T\otimes E\otimes L_0^m),\,d\bigr)$。特に $k>n$ なら $\HH^k=0$、
\item[(c)] $\HH^0(E\otimes L_0^m)=$ 平坦（Casimir）正則切断全体。
\end{enumerate}
\end{thm}
\begin{hyptab}
$X$ 射影的、$L_0$ 豊富 & 古典 Serre の仮定そのもの。 & (a)：Hartshorne III.5.2 を有限個（$b=0,\dots,n$）の連接層 $\wedge^bT\otimes E$ に適用する。\\\midrule
$L_0$ が Poisson 加群（$E$ も） & $L_0^m$ にテンソル積の加群構造が入り、$C(A;E\otimes L_0^m)$ が複体になる。 & 複体 $K^\bullet=C(A;E\otimes L_0^m)$ の定義。正値性とは無関係。\\\midrule
$[\sigma,\sigma]=0$ & $d^2=0$。 & 同上。\\\midrule
（不要）$\sigma$ の非退化性・正規交差 & --- & この定理では使わない。
\end{hyptab}
\begin{proof}
$L_0^m$ はテンソル積の加群構造（Leibniz 則）をもつ。Serre の定理（定理 \ref{thm:Serre}、Hartshorne III.5.2）を $n+1$ 個の連接層 $\wedge^bT\otimes E$ に適用し、$m_0$ をそれらの最大にとれば (a)。各 $K^b$ は $\Gamma$-非輪状なので、有界な複体 $K^\bullet$（微分は $\C$-線形で $\mathcal O$-線形でない）に対して、スペクトル系列 $E_1^{a,b}=H^a(K^b)$ は $a=0$ に集中し、$R\Gamma(K^\bullet)=\Gamma(K^\bullet)$（非輪状分解の補題）。これが (b)。$b\le n$ なので $k>n$ で消える。(c) は $\Gamma$-複体の $H^0$ で、$\HH^0=\ker(d_A:H^0(E)\to H^0(T\otimes E))$、すなわち平坦切断。\qedhere
\end{proof}
\begin{rem}[全次数での消滅は偽]\tP
$\sigma=0$（複素構造）の場合、$\HH^k(L)=\bigoplus_{a+b=k}H^a(X,\wedge^bT\otimes L)$ なので、$m\gg0$ でも $H^0(\wedge^bT\otimes L^m)\ne0$ により $\HH^b\ne0\ (b=0,\dots,n)$。$\mathbb{CP}^2$ の $m=0$ では $[1,8,10,0,0]$（\tV）。消えるのは横断次数 $a>0$ と全次数 $>n$ のみで、「全次数での Serre 消滅」は成り立たない。
\end{rem}
\begin{rem}[\tV\ 有効な閾値]
Bott の公式から、$E_1^{a>0}=0$ となる条件は $\mathbb{CP}^2$ で $m\ge-2$（$m=-3$ では $(2,0),(1,1),(0,2)$ の3項がすべて非零）、$\mathbb{CP}^1\times\mathbb{CP}^1$（$L_0=\mathcal O(1,1)$）で $m\ge-1$（$m=-2$ は $K_X$ で例外）。どちらも、計算した表で次数 $\ge3$ が消える範囲と一致する。
\end{rem}

\subsection{周期的安定化}
\begin{thm}[周期的安定化（対数シンプレクティックな曲面）]\tP\ （定理 \ref{thm:main2} とその仮定に依存）\label{thm:A2}
$X$ を単連結な射影曲面、$\sigma$ を $D$（反標準、通常二重点のみ）に沿って対数シンプレクティック、$L_0$ を Poisson 加群をもつ直線束で留数が $\rho_i=-a_i$（$a_i\in\mathbb Q_{>0}$、各成分に沿って定数）、$U=X\setminus D$ をアファイン（$D$ 豊富）とする。$N$ を $a_i$ の分母の最小公倍数とする。$m\ge1$ ならば
\[\HH^k(L_0^m)=0\ (k\ge3),\qquad\HH^*(L_0^m)=H^*(U;L_U^{(m)})\oplus\C^\mu[-2],\]
ここで $L_U^{(m)}$ は $D_i$ まわりのモノドロミー $e^{-2\pi ima_i}$ の局所系。特に $\HH^*(L_0^m)$ は $m\bmod N$ のみに依存する。
\end{thm}
\begin{hyptab}
留数 $-ma_i<0$ & $r_i\ne1$ かつ正整数でない。 & 定理 \ref{thm:main2}(i)：$\mathcal F=Rj_*L_U$。\\\midrule
$U$ アファイン（$D$ 豊富） & Artin 消滅 $H^k(U;L)=0\ (k>2)$。 & 次数 $\ge3$ の消滅。（定理 \ref{thm:A1} (b) は射影性と豊富性で同じ結論を与え、アファイン性を不要にする。）\\\midrule
$X$ 単連結 & $\pi_1(U)$ のアーベル化が子午線で生成される。 & 局所系が $m\bmod N$ のみで決まる。\\\midrule
$a_i\in\mathbb Q_{>0}$ & $c_1(L_0)$ が境界成分の正の有理結合（「境界に関して正」）。 & 周期性。
\end{hyptab}
\tV\ $\mathbb{CP}^2$（$L_0=\mathcal O(1)$、$a=1/3$）で周期 3、$\mathbb{CP}^1\times\mathbb{CP}^1$（$L_0=\mathcal O(1,1)$、$a=1/2$）で周期 2、$m\ge1$。計算した8例すべてで確認。

\begin{conj}[一般の場合]\tK\label{conj:A3}
一般の一般化複素多様体 $M$ と、適切な意味で「正」な一般化正則直線束 $L_0$（候補：正則点の葉空間上で豊富で、型が跳ぶ集合に沿った留数が正）に対し、$\HH^*(L_0^m)$ は $m\gg0$ で全次数 $>n_{\max}$ で消え、$m$ について周期的になる。$\beta$-変換や正則型でない一般の $\Jc$ では、「連接な濾過 $K^b$ の $E_1$ 項」を定義する方法がなく、これが定理 \ref{thm:A1} を超える際の主な障害である。
\end{conj}

\subsection{Pal--Poddar の定理との関係}
\begin{prop}[積の場合の整合性]\tP\label{prop:PPprod}
$M=Y\times X$（$Y$ コンパクトシンプレクティック、$X$ 射影的複素多様体）、$E=\mathrm{pr}_Y^*F\boxtimes\mathrm{pr}_X^*G_0^m$（$F$ は $Y$ 上の平坦束、$G_0$ は豊富）とする。このとき、一般化 Dolbeault 複体は $Y$ の de Rham 複体と $X$ の複体のテンソル積で、
\[\HH^k=\bigoplus_{j+b+a=k}H^j_{dR}(Y;F)\otimes H^a(X,\wedge^bT_X\otimes G_0^m)\]
（Künneth）。古典 Serre により $m\gg0$ で $a>0$ の項は消え、残る $H^j_{dR}(Y;F)\otimes H^0(X,\wedge^bT_X\otimes G_0^m)$ は一般に全次数 $\dim Y+\dim_\C X$ まで非零。したがって、$\mathrm{pr}_X^*G_0^m$ に対する Serre 型消滅は横断次数 $a$ についてのみ成り立つ。
\end{prop}
\src{Pal--Poddar, \emph{Strong generalized holomorphic principal bundles}, arXiv:2404.18113, \S11, Thm.\ 11.12：$M$ がコンパクトな正則型（型 $k$）一般化複素多様体で $M/\mathcal S$ が Kähler 軌道体、$E$ が正の SGH 直線束のとき、(1) $p+q>k$ ならば $H^q(M,(\mathcal G^*M)^p\otimes E)=0$、(2) 任意の SGH ベクトル束 $E'$ に対し $m_0$ が存在して $m\ge m_0$、$q>0$ ならば $H^q(M,E'\otimes E^m)=0$。}
\begin{rem}[\tQ]
Pal--Poddar の $H^q(M,(\mathcal G^*M)^p\otimes E)$（Cor.\ 8.10 により $H^{p,q}_{d_L}(M,E)$）は、葉に沿って定数な横断形式の層コホモロジーで、$q$ は横断（複素）次数である（\S\S5--11, 13 を再読して確認。抽出は要約を介しており式の細部は未確認）。したがって上の命題の読み（横断次数 $a$ についてのみ消滅）は Thm.\ 11.12 と整合する。ただし群は同じではない：Pal--Poddar は $\mathcal G^*M$ の外冪、本ノートは $\wedge^bT$（接側）をテンソルし、さらに本ノートの $\HH$ は葉方向の de Rham 因子を含む全複体である。上の命題は積の場合の計算であり、Thm.\ 11.12 そのものの検証ではない。定理 \ref{thm:A1} は正則型や葉空間の軌道体を仮定せず、型が跳ぶ $\beta$-変換（例：型 0 と型 $n$ が共存）も含む代わりに、基礎の複素多様体が射影的であること（$\beta$-変換の場合に限る）を必要とする。
\end{rem}

%--------------------------------------------------------------
\section{双対性}\label{sec:duality}
%--------------------------------------------------------------
\begin{thm}[一般化正則直線束の Serre 双対性]\tP\label{thm:dual}
$X$ を任意のコンパクト複素多様体（$\dim_\C X=n$）、$\sigma$ を任意の正則 Poisson 構造（Kähler、射影、対数シンプレクティック、非退化のいずれも不要。$\sigma=0$ を含む）、$L$ を Poisson 加群である直線束とすると、すべての $k$ で
\[\HH^k(L)^*\;\cong\;\HH^{2n-k}\bigl(L^\vee\otimes K_X^{2}\bigr).\]
$L^\vee\otimes K_X^2$ はテンソル積の加群構造をもつ（枠 $e'$、$ee'=\omega^2$ で $V'=-V+2V_K=-V-2M$。$K_X$ の標準加群構造は $V_K^k=\sum_i\partial_i\sigma^{ki}$、$M^j=\sum_i\partial_i\sigma^{ij}$ は $\omega=dz_1\wedge\cdots\wedge dz_n$ に関する modular ベクトル場）。
\end{thm}
\begin{hyptab}
$X$ コンパクト & 楕円型複体の閉値域性・有限次元性、Stokes で境界項が出ない。 & 手順 (4)（双対の同定）。\\\midrule
$[\sigma,\sigma]=0$、$L$ が Poisson 加群 & $D^2=0$、$D'^2=0$（$L'$ も加群）。 & 複体として扱うため。\\\midrule
（不要）Kähler / 非退化 / 対数 & 楕円性を主シンボルから直接確かめ、転置は局所的な作用素の等式なので大域的な原始関数が不要。 & 手順 (1),(3)。
\end{hyptab}
\begin{proof}[証明（概略）]
(0) Dolbeault モデル：$\mathrm{Tot}^k(L)=\bigoplus_{a+b=k}A^{0,a}(\wedge^bT\otimes L)$、$D=\dbar+d_K$（$d_K(\phi\otimes Pe)=(-1)^a\phi\otimes d(Pe)$、$d=[\sigma,\cdot]+V\wedge$）。$D^2=0$ は $\dbar^2=0$、$d^2=0$、$\dbar d_K+d_K\dbar=0$ による（命題 \ref{prop:ident}(b)）。
(1) 楕円性（\tP）：実余接ベクトル $\xi\ne0$ での $D$ の主シンボルは $v(\xi)=(\xi^{0,1},\sigma^\sh\xi^{1,0})\in L^*=T^{*0,1}\oplus T^{1,0}$ との外積で、$(\wedge L^*,v\wedge)$ は $v\ne0$ のときに限り完全。$v=0$ なら $\xi^{0,1}=0$、$\xi$ 実なので $\xi=0$。よって $\sigma$ が退化する点も含め、すべての点で楕円型（Gualtieri の命題 \ref{prop:ellip} をこの場合に直接確かめたもの）。$L'$ についても同様。
(2) 完全対：$E^b(L)\times E^{n-b}(L')\to\wedge^nT\otimes K^2=K$、$A^{0,a}\times A^{0,n-a}\to A^{0,n}(K)=A^{n,n}$ を積分して $\mathrm{Tot}^k(L)\times\mathrm{Tot}^{2n-k}(L')\to\C$（各点で完全対）。
(3) 転置：次の2つの恒等式から、形式的転置 $D^t$ が $-(-1)^{\deg}D'$ になる。
\end{proof}
\begin{lem}[補題 D：Schouten 括弧の最高次]\tP\ \tV\label{lem:D}
$T=\sum_j(-1)^ju_j\,\partial_1\wedge\cdots\widehat{\partial_j}\cdots\wedge\partial_n\in\wedge^{n-1}T$ に対し $[\sigma,T]_{\rm top}=-\sum_{i,j}\partial_i(\sigma^{ij}u_j)+2\sum_jM^ju_j$、$(M\wedge T)_{\rm top}=\sum_jM^ju_j$。
\end{lem}
\begin{lem}[補題 E：転置恒等式]\tP\ \tV\label{lem:E}
$P\in\wedge^bT$、$Q\in\wedge^{n-b-1}T$（滑らかな係数）、$d=[\sigma,\cdot]+V\wedge$、$d'=[\sigma,\cdot]+V'\wedge$、$V'=-V-2M$ とするとき、$\langle P,Q\rangle:=(P\wedge Q)_{\rm top}$ について
\[\langle dP,Q\rangle+(-1)^b\langle P,d'Q\rangle=-\sum_i\partial_i\Bigl(\sum_j\sigma^{ij}u_j\Bigr),\quad T=P\wedge Q=\sum_j(-1)^ju_j\partial_{\hat j}.\]
証明：Leibniz 則 $[\sigma,P\wedge Q]=[\sigma,P]\wedge Q+(-1)^bP\wedge[\sigma,Q]$（$\sigma$ は偶数次数）、$V$-項は $V\wedge P\wedge Q+(-1)^bP\wedge(-V)\wedge Q=0$、$M$-項は $-2M\wedge T$。補題 D から従う。
\tV\ $n=2,3,4$ で、ランダムな多項式・Laurent 単項式データ、任意の二重ベクトル場 $\sigma$（Poisson でなくても恒等式自体は成立）に対し、恒等式を厳密算術で確認（\texttt{cp2/duality\_ndim\_exact.py}）。定数 $c=-2$ と符号 $(-1)^b$ は、導出とは独立に連立方程式として解いて求めた（\texttt{cp2/duality\_ndim.py}）。$[\sigma,M]=0$ と $d^2=0$（$V=M$、$\sigma$ Poisson）も確認。
\end{lem}
\begin{proof}[定理 \ref{thm:dual} の証明の続き]
(3) 続き：補題 E の右辺は $\partial_iF^i\,\omega\wedge d\bar z^{1\cdots n}=\pm d(F^i\iota_{\partial_i}\omega\wedge d\bar z^{1\cdots n})$ の形の全微分（$d\bar z^{1\cdots n}$ との積で $\dbar$ 成分が消える）で、Stokes により積分は 0。$\dbar$ 部分は古典的な Serre 双対性の計算で $\langle\dbar\alpha,\beta\rangle=-(-1)^{a+b}\langle\alpha,\dbar\beta\rangle$。$d$ 部分は $\langle d_K\alpha,\beta\rangle=-(-1)^{a+b}\langle\alpha,d'_K\beta\rangle$。2つの符号が一致することを手計算で確認した（機械では未確認\tQ）。$D^t$ は局所作用素なので、座標近傍での等式から大域的に $\langle D\alpha,\beta\rangle=-(-1)^{\deg\alpha}\langle\alpha,D'\beta\rangle$。
(4) 結論：$D,D'$ は閉多様体上の楕円型複体で、完全対に関して互いの転置（符号を除く）。楕円型複体の Hodge 理論により $\HH^k$ は有限次元で値域は閉、$\HH^k(L)^*$ は転置複体のコホモロジー（超関数係数）で、楕円正則性により滑らかな係数のもの $\HH^{2n-k}(L')$ に等しい。\qedhere
\end{proof}
\begin{rem}[整合性の確認\ \tV]
(i) $\sigma=0$：古典的 Serre 双対性（$\wedge^bT^*=\wedge^{n-b}T\otimes K$ を使う）。(ii) $\mathbb{CP}^2$、$K^2=\mathcal O(-6)$：$m\leftrightarrow-m-6$、$\mathbb{CP}^1\times\mathbb{CP}^1$、$K^2=\mathcal O(-4,-4)$：$m\leftrightarrow-m-4$（\S\ref{sec:cp2}, \S\ref{sec:main2} の表と一致）。(iii) 対数シンプレクティックな場合、$K_X$ の留数は 1（命題 \ref{prop:K} の $r=1$）なので $L'$ の留数は $2-r_i$、これが主定理の公式の $r\leftrightarrow2-r$ の対称性の起源（任意の偶数次元で）。
\end{rem}

%--------------------------------------------------------------
\section{Riemann--Roch 型の公式}\label{sec:RR}
%--------------------------------------------------------------
\subsection{全 Euler 標数は位相不変量}
\begin{prop}\tP\label{prop:euler}
$X$ をコンパクト複素多様体、$\sigma$ を任意の正則 Poisson 構造、$E$ を階数 $\rk E$ の Poisson 加群とする。このとき
\[\sum_k(-1)^k\dim\HH^k(E)=\sum_b(-1)^b\chi(X,\wedge^bT\otimes E)=(-1)^n\,\rk E\cdot\chi(X).\]
\end{prop}
\begin{proof}
複体 $C(A;E)=(\wedge^bT\otimes E,d)$ は $b$ による濾過（stupid filtration）をもち、スペクトル系列 $E_1^{a,b}=H^a(X,\wedge^bT\otimes E)$ は有限次元で $\HH^{a+b}$ に収束する。Euler 標数は $E_1$ 項のものに等しく、$\sum_b(-1)^b\chi(\wedge^bT\otimes E)$ は Hirzebruch--Riemann--Roch（定理 \ref{thm:HRR}）で
$\int\ch(E)\Td(T)\sum_b(-1)^b\ch(\wedge^bT)$。$\Td(T)=\prod\frac{x_i}{1-e^{-x_i}}$、$\sum_b(-1)^b\ch(\wedge^bT)=\prod(1-e^{x_i})$ なので積は $\prod(-x_ie^{x_i})=(-1)^ne^{c_1}c_n(T)$。$c_n$ が最高次なので $\ch(E)$ からは階数のみが効き、$(-1)^n\rk E\int c_n(T)=(-1)^n\rk E\,\chi(X)$。\qedhere
\end{proof}
\begin{rem}[\tP\ 意味するところ]
$\sigma=0$ のときこの公式は $\mathbb{P}^1$、$E=\mathcal O(m)$ で $\chi(\mathcal O(m))-\chi(\mathcal O(m+2))=(m+1)-(m+3)=-2$ と確認できる（$m$ に依存しない）。つまり、一般化 Dolbeault 複体の「全」Euler 標数に対する Riemann--Roch 公式は、$c_1(E)$ を見ない位相的な公式であり、古典の $\chi(M,E)=\int\ch(E)\Td(M)$ の一般化にはなっていない。古典の $\chi(M,E)$ は $b=0$ の「列」$H^\bullet(X,E)$ の Euler 標数にあたる。一般の $\Jc$ ではこの列は $d$ で保たれず、何らかの細分が必要になる。
\end{rem}

\subsection{Poincaré 多項式の閉公式}
\begin{thm}[Riemann--Roch 型の閉公式（曲面）]\tP\ （定理 \ref{thm:main2} と下の仮定に依存）\label{thm:B}
定理 \ref{thm:main2} の設定で、次を追加で仮定する。(H1) $H^1(X;\C)=0$（$b_1(U)=c-\rho$ のため）、(H2) $U$ がアファイン（次数 $\ge3$ の消滅のため）、(H3) 全成分で留数が等しい $r_i=r\ne1$、(H4) $r\notin\Z$ のとき、$H^0(U;L_U)=H^1(U;L_U)=0$（G2：下の注意）。$b_1=c-\rho$（$c$ は成分数、$\rho$ は $[D_i]\in H^2(X;\C)$ の張る空間の次元）、$\chi_U=\chi(X)-\mu$、$b_2=\chi_U-1+b_1$ とおく。このとき
\begin{align*}
P_E(t)&=1+b_1t+(b_2+\mu)t^2 &&(r\in\Z_{\le0}),\\
P_E(t)&=(b_2+\mu)t^2+b_1t^3+t^4 &&(r\in\Z_{\ge2}),\\
P_E(t)&=\chi(X)\,t^2 &&(r\notin\Z).
\end{align*}
同値な書き方：$P_E(t)=\chi(X)t^2+[r\le0]\,\Delta(t)+[r\ge2]\,t^4\Delta(1/t)$、$\Delta(t)=(1+t)(1+(b_1-1)t)$。
\end{thm}
\begin{hyptab}
(H1) $H^1(X;\C)=0$ & Gysin 完全列 $H^0(D)\to H^2(X)$ から $H^1(U)=\ker(\C^c\to H^2(X))$。 & $b_1(U)=c-\rho$。\\\midrule
(H2) $U$ アファイン & Artin 消滅。 & $b_3=b_4=0$、したがって $b_2=\chi_U-1+b_1$。\\\midrule
(H3) 留数が等しい & $\mathcal F$ が $Rj_*$ か $j_!$ のどちらか一方になる。 & 定理 \ref{thm:main2}(i),(ii)。\\\midrule
(H4) $r\notin\Z$ で $H^0(U;L)=H^1(U;L)=0$ & $H^*(U;L)=[0,0,\chi_U]$ になるため。 & $r\notin\Z$ の行。$\mathbb{CP}^2$ の結節曲線では $\pi_1(U)$ がアーベル群（Fulton, Ann.\ Math.\ 111 (1980) 407--409；Deligne, Sém.\ Bourbaki 543）。\tC\\\midrule
$r_i\ne1$ & 定理 \ref{thm:main2} の仮定。 & 補題 \ref{lem:Q}。
\end{hyptab}
\begin{proof}
定理 \ref{thm:main2}(i),(ii) に $H^*(U)=[1,b_1,b_2,0,0]$（$r\in\Z_{\le0}$、自明な局所系）、$H_c^*(U)$（Poincaré 双対で $[0,0,b_2,b_1,1]$、$r\in\Z_{\ge2}$）、$H^*(U;L)=[0,0,\chi_U,0,0]$（$r\notin\Z$：(H2)、(H4)、$\chi(U;L)=\chi(U)$ から）を代入し、$\mu t^2$ を加える。同値な形は $\chi(X)=\chi_U+\mu=1-b_1+b_2+\mu$ による。
\end{proof}
\begin{rem}[性質\ \tP]
(i) $P_E(-1)=\chi(X)$、すべての $r$ で（命題 \ref{prop:euler}、$n=2$ と整合）。(ii) 双対性 $P_{2-r}(t)=t^4P_r(1/t)$ は式から明らか（定理 \ref{thm:dual} の公式レベルの帰結）。(iii) $\mu$ と $b_2$ は $\chi(X)t^2$ に吸収され、跳びの項 $\Delta$ は $b_1$（$[D_i]$ の間の線形関係の数）のみに依存する。(iv) $E=L_0^m$ のとき、古典の Riemann--Roch が与える $m$ の擬多項式は定数 $\chi(X)$ で、$m$ への依存は留数 $r=-ma$ の整数性（周期 $N$）と符号のみに現れる。
\end{rem}
\begin{rem}[\tQ\ G2 について]
(H4) を、検証した8例（$\mathbb{CP}^2$ の結節曲線 4 種、$\mathbb{CP}^1\times\mathbb{CP}^1$ の 4 種）で確かめた。$\mathbb{CP}^1\times\mathbb{CP}^1$ の滑らかな $(2,2)$ 曲線と 2 本の $(1,1)$ 曲線については、Timmerscheidt（\emph{Mixed Hodge theory for unitary local systems}, J.\ reine angew.\ Math.\ 379 (1987) 152--171）と Arapura（\emph{Geometry of cohomology support loci for local systems I}, arXiv:alg-geom/9612006、J.\ Algebraic Geom.\ 6 (1997) 563--597（巻頁は記憶\tQ）、トーション平行移動の定理）を引用した（適用範囲の確認は未了\tQ）。一般には $H^1(U;L_\chi)=0\Leftrightarrow\chi\notin V^1(U)$（第 1 特性多様体）で、非ノーダルな $D$（例：3 本共点線）では成り立たないことがある（Libgober, Arapura）。
\end{rem}
\tV\ 閉公式は、$\mathbb{CP}^2$（滑らか・結節・円錐＋割線・三角形、$m=-15..15$）と $\mathbb{CP}^1\times\mathbb{CP}^1$（4 例、$m=-9..8$）の計 196 個の値で厳密計算と一致（食い違い 0、\texttt{cp2/ncd\_global\_check.py}）。双対性 $r\leftrightarrow2-r$ の閉公式レベルの対称性は 73 個の有理数 $r$ の格子で確認（\texttt{cp2/serre\_rr\_check.py}）。

\subsection{指数理論的な見通し}\label{sec:index}
本節の内容は予想であり、指数定理は何も証明していない。
\begin{obs}[\tK\ 指標根の描像（定性的）]
滑らかな点で対数複体は単項式 $z^ne$ に $(n-r)$ 倍で作用する（補題 \ref{lem:stalk}）。指標根（indicial root）は $r\in\Z$ のときに限り存在し、$r\notin\Z$ ではウェイトごとの複体が非輪状で $P_r=\chi(X)t^2$。$r\in\Z$ で残るのは、$D$ の管状近傍の境界 $\partial U$ のコホモロジー（Gysin 完全列で $H^*(D)\otimes(1\oplus t)$）の部分商で、$\chi(\partial U)=0$ なので Euler 標数には見えない。これは Melrose の APS／b-計算の描像（Melrose, \emph{The Atiyah--Patodi--Singer Index Theorem}, A K Peters 1993）と定性的に整合する。
\end{obs}
\begin{conj}[指数型 Riemann--Roch（対数シンプレクティック曲面）]\tK\label{conj:I3}
(a) $\sum(-1)^k\dim\HH^k(L)=\int_Xc_2(T_X(-\log D))+\mu=\chi(X)$ は、境界項なしの Atiyah--Singer 型積分で与えられる。(b) 跳びの項 $J(r)=P_r(t)-\chi(X)t^2$ は $r\in\Z\setminus\{1\}$ でのみ非零で、$r\le0$ で $\Delta(t)$、$r\ge2$ で $t^4\Delta(1/t)$。$\Delta(t)$ は指標根の核 $H^*(\partial U)$ のある部分商の Poincaré 多項式（長完全列 $H_c^*(U)\to H^*(U)\to H^*(\partial U)$ の重さ濾過で固定される）。(c) 一般次元・非正規交差の $D$ では、跳びの項は特異点のスペクトル数で支配される局所項（\S\ref{sec:local}）を含む。
\end{conj}
\begin{rem}[\tC\ 文献の状況と障害]
\begin{itemize}
\item \emph{$b$-指数との類似}：Braverman--Loizides--Song, \emph{Geometric quantization of $b$-symplectic manifolds}, arXiv:1910.10016 は、$b$-シンプレクティック多様体の量子化指数が $\int\mathrm{Td}({}^bTM)\mathrm{Ch}(L)$ で与えられ、境界項が極限で消えることを示す（題名・著者は確認した\tC）。本ノートの Euler 標数が $\chi(X)=\int c_2(T_X(-\log D))+\mu$ 型で $r$ に依らないことと類似する。
\item \emph{Lie 亜群の指数}：Pflaum--Posthuma--Tang, \emph{The index of geometric operators on Lie groupoids}, arXiv:1308.0236（Indag.\ Math.\ 25 (2014) 1--18）；Qiao--So, arXiv:2304.13345（境界亜群、余次元 $\ge2$ で eta 項なし）；Ammann--Lauter--Nistor, arXiv:math/0304044；Mazzeo--Melrose, arXiv:math/9812120；Hausel--Hunsicker--Mazzeo, arXiv:math/0207169（Duke Math.\ J.\ 122 (2004)、頁は\tQ）。題名・著者は確認したが、本問題への適用可能性は未確認。
\item \emph{Poisson コホモロジー}：Lanius, arXiv:1605.03854（対数シンプレクティック多様体の $C^\infty$ Poisson コホモロジー）。本ノートの正則 Poisson 超コホモロジーとは別の対象で、「$b$-部分 $+$ 境界項」の構造の類似があるのみ。
\item \emph{見つからなかったもの}：複素 Lie 亜群 $L=T^{0,1}\bowtie A^{1,0}$ に対する指数定理、留数を重みとする $b$-Dolbeault 指数。実 Lie 亜群 $T_X(-\log D)$ は Gualtieri--Li, arXiv:1206.3674 により滑らかな $D$ で積分可能。\tC\ 正規交差の場合は未確認。
\item \emph{単模性}：Pflaum--Posthuma--Tang のトレースには不変密度が必要で、Poisson 構造では modular 類の消滅にあたる。対数シンプレクティックでは対数体積形式から modular ベクトル場が 0 になる見込みだが未確認\tQ。
\end{itemize}
\end{rem}

%--------------------------------------------------------------
\section{トーリックな場合と高次元}\label{sec:toric}
%--------------------------------------------------------------
\subsection{トーリック対数シンプレクティック構造}
$X$ を滑らかな射影的トーリック多様体（$N=\Z^n$、$n=2k$）、光線 $u_i$ ($i\in I$)、$D$ をトーリック境界（単純正規交差）とする。$E_j=z_j\partial_j$（対数枠）に関し、定数の非退化歪対称行列 $A$ により $\sigma=\tfrac12\sum A_{ij}E_i\wedge E_j$（トーラス不変、$D$ に沿って対数シンプレクティック）。$L=\mathcal O(\sum a_iD_i)$ に平坦対数接続 $d+\rho$（$\rho\in M_\C$）を与える。留数は $r_i=\langle\rho,u_i\rangle-a_i$、$c_1(L)=-\sum r_i[D_i]$。\tP\ すべての加群はこの形（$H^0(\Omega^1(\log D))=M_\C$、補題 \ref{lem:res}）。

\begin{thm}[weight 公式]\tP\ \tV\label{thm:weight}
$E_1^{a,b}=H^a(X,\wedge^bT\otimes L)=0\ (a>0)$ を仮定する（例：定理 \ref{thm:A1} の範囲）。$\chi\in M$ に対し $c_i(\chi)=\langle\chi,u_i\rangle+a_i$、$C(\chi)=\{i:c_i(\chi)=-1\}$、$W(\chi)=\mathrm{span}(u_i:i\in C(\chi))$、$s=\dim W$ とおく。このとき
\[P_{\HH}(t)=\sum_{\chi\in\Lambda}t^{s(\chi)}(1+t)^{n-s(\chi)},\quad\Lambda=\{\chi\in M:c_i(\chi)\ge-1\ \forall i,\ A(\chi+\rho)\in W(\chi)\}.\]
\end{thm}
\begin{hyptab}
トーリック・トーラス不変な $\sigma$、定数 $A$ & 大域切断の複体がトーラスの指標 $\chi\in M$ で次数付き（$T$-weight）に分解され、各 weight で有限次元の線形代数になる。 & 証明：weight $\chi$ の部分複体 $G^b_\chi$ の決定。\\\midrule
$E_1^{a>0}=0$ & $\HH^*$ が大域切断の複体のコホモロジーに等しい。 & 定理 \ref{thm:A1}(b)。\\\midrule
非退化な $A$ & $\sigma^\sh:B\cong T_D$ が成立し、$\wedge^b\sigma^\sh$ で対数 de Rham 微分 $(\chi+\rho)\wedge$ が $v_\chi\wedge$（$v_\chi=A(\chi+\rho)$）に写る。 & 補題 \ref{lem:iso}。
\end{hyptab}
\begin{proof}[証明の概略]
$z^\chi w$（$w\in\wedge^bN_\C$）が $\wedge^bT\otimes L$ の正則切断であるための条件は、各 $D_i$ に沿って $c_i(\chi)\ge-1$、かつ $c_i(\chi)=-1$ なら $u_i\wedge w=0$。したがって weight $\chi$ の部分は $G^b_\chi=\{w\in\wedge^bN_\C:u_i\wedge w=0\ (i\in C(\chi))\}$（ある $c_i<-1$ なら空）で、$G^\bullet_\chi\cong\wedge^{\bullet-s}(N/W)\otimes\wedge^sW$ 型。微分は $v_\chi\wedge$。$v_\chi\bmod W=0$ のときは微分 0 で Poincaré 多項式 $t^s(1+t)^{n-s}$、そうでなければ非輪状。
\end{proof}
\tV\ 閉公式と、線形代数による直接計算（\texttt{cp2/toric\_hh.py}）が 60 個のランダムな場合（$\mathbb P^2$、$\mathbb P^1\times\mathbb P^1$、$\mathbb P^4$、$\mathbb P^2\times\mathbb P^2$、$(\mathbb P^1)^4$）で一致。

\begin{cor}\tP
(a) $E_1^{a>0}=0$ なら全 Euler 標数は $s(\chi)=n$ となる weight の個数（頂点の個数）$=\chi(X)$。（命題 \ref{prop:euler} と整合。）(b) 寄与の由来：$s=0$ は $A(\chi+\rho)=0$、すなわち $\chi=-\rho$ で、$\rho\in M$ かつ全 $r_i\le0$ のとき $(1+t)^n=H^*(U)$（$Rj_*$ の場合）。$s=n$ は常に条件を満たし頂点 weight、$t^n$ を各 1 個。$0<s<n$ は $x=\chi+\rho\in A^{-1}W$ かつ $\langle x,u_i\rangle=r_i-1\ (i\in C)$ という条件が必要で、積分格子上の「共鳴」でのみ解がある。(c) $s=n-1$ で全 $r_i=1\ (i\in C)$ のとき、条件は自動的に成り立ち、トーリック曲線 $C$ 上の格子点の個数 $\deg(L|_C)+1$ 個のクラス $t^{n-1}(1+t)$ が現れる（次元 2 での $r_i=1$ の共鳴に対応）。
\end{cor}

\subsection{一般の $\sigma$ と 4 次元}
\begin{obs}[\tV]
$A$ が一般的（成分が非常に大きい整数の組）のとき、$\Lambda=\{\text{頂点}\}\cup\{\text{共鳴}\}\cup\{-\rho\}$（$-\rho$ は $\rho\in M$ かつ全 $r_i\le1$ のとき含まれる）となり、$P_{\HH}=$（$U$-部分）$+\#\text{頂点}\cdot t^n+$（$r=1$ の共鳴クラス）。$U$-部分は $\rho\in M$ かつ全 $r_i\le0$ のとき $(1+t)^n$、そうでなければ 0。素数 $2^{61}-1$ を法とする厳密算術で 3626 個の場合（$\mathbb P^4$ ($m=1,2$)、$\mathbb P^2\times\mathbb P^2$、$(\mathbb P^1)^4$、$\mathbb P^1\times\mathbb P^3$、$\mathbb P^1\times\mathbb P^1\times\mathbb P^2$、$\rho\in\{-2,..,1\}^4$ と 3 つの非整数 $\rho$）で食い違い 0。例：$\mathbb P^4$、$m=1$：$[1,4,6,4,6]$；$\mathbb P^2\times\mathbb P^2$：$[1,4,6,4,10]$。
\end{obs}
\begin{obs}[新しい現象：$\sigma$ の算術への依存\ \tV]\label{obs:arith}
次元 4 では $\HH^*$ は $(X,D,L,\text{留数})$ だけでは決まらない。$\mathbb P^2\times\mathbb P^2$ の積構造 $\sigma$（ブロック対角 $A$）では $[1,4,12,16,16]$（2 つの曲面の結果の積、Künneth）、一般の $A$ では $[1,4,6,4,10]$。差 $(0,0,6,12,6)=6t^2(1+t)^2$ は、余次元 2 の 6 個の層 $\{\text{一方の因子の二重点}\}\times(\mathbb C^*)^2$ の $t^2H^*(O_\tau)$。仕組み：$|C|=2$ の層で、$u_i,u_j$ が $\sigma$ に関してシンプレクティックなら $x_C(A,r)\in A^{-1}W$ が一意に決まり、これが $\rho$ の格子平行移動になる算術的条件のときだけ寄与する。\tK\ $\HH$ は対数シンプレクティック構造のパラメータ空間の可算個の真解析部分集合上で上に跳び（一般化 Dolbeault コホモロジーの上半連続性）、一般の値は最小である。
\end{obs}
\begin{conj}[一般の単純正規交差 $D$]\tK
一般の（トーリックでない）単純正規交差の対数シンプレクティック $(X^{2k},D)$、一般の $\sigma$、$r_i\ne1$ の定数留数の $L$ に対し、$\HH^*(L)=H^*(X;\mathcal F)\oplus\bigoplus_{p\in D_{(n)}}\C[-n]$（$D_{(n)}$ は $n$ 本の成分が交わる点）。Euler 標数との整合は、低次の層 $D_k^\circ$ の $\chi$ が消える場合に限る（トーリックでは成立、一般には未証明）。上の \S\ref{sec:toric} の計算はトーリックのみ。
\end{conj}

%--------------------------------------------------------------
\section{非ノーダルな特異点の局所理論}\label{sec:local}
%--------------------------------------------------------------
$D$ に尖点・tacnode・3重点などがあるとき、定理 \ref{thm:main2} は使えない。以下は計算による観察と予想である。局所モデル：$\sigma=h\,\partial_z\wedge\partial_w$、加群 $\mathcal O\,h^s$、$d_0f=(hf_w+sfh_w)\partial_z-(hf_z+sfh_z)\partial_w$、$d_1(p,q)=h(p_z+q_w)-(1-s)(h_zp+h_wq)$。大域では $s=r=-m/3$（$\mathbb{CP}^2$）、$-m/2$（$\mathbb{CP}^1\times\mathbb{CP}^1$）。

\subsection{茎のコホモロジー}
\tV\ 準斉次な孤立特異点の germ について、重み付き厳密線形代数（スクリプト \texttt{local\_stalk.py} など）で茎の次元 $(h^0,h^1,h^2)$ を計算した（形式的 germ のコホモロジーが収束 germ のものに等しいことを仮定）：
\begin{center}\small
\begin{tabular}{lccc}\toprule
特異点 & 一般の $s$ & $s$ 整数 $\le0$ & $s$ がスペクトル数 mod $\Z$、$s<\mathrm{lct}$\\\midrule
滑らか ($z$) & $(0,0,0)$ & $(1,1,0)$ & ---\\
二重点 ($zw$) & $(0,0,1)$ & $(1,2,2)$ & ---\\
尖点 ($z^2-w^3$) & $(0,0,2)$ & $(1,1,2)$ & $(0,1,3)$\\
tacnode ($z^2-w^4$) & $(0,0,3)$ & $(1,2,4)$ & $(0,1,4)$\\
$A_4$ ($z^2-w^5$) & $(0,0,4)$ & $(1,1,4)$ & ---\\
$D_4$ ($zw(z-w)$) & $(0,0,4)$ & $(1,3,6)$ & $(0,1,5)$\\
$D_5$ ($w(z^2-w^3)$) & $(0,0,5)$ & $(1,2,6)$ & $(0,1,6)$\\\bottomrule
\end{tabular}
\end{center}
\begin{obs}[\tV\ 経験則]
スペクトル数（lct は最小スペクトル数：滑らか・二重点 1、尖点 $5/6$、tacnode $3/4$、$D_4$ $2/3$、$D_5$ $5/8$）に関し、$s\ne1$ で
\[\text{茎}_p=\C^{\mu_p}[-2]\ \oplus\ \begin{cases}\text{茎}(Rj_*L_s)_p&(s<\mathrm{lct}_p)\\0&(s>\mathrm{lct}_p)\end{cases}\]
（$\mu_p$ は Milnor 数 $=$ Tjurina 数）。$Rj_*$ の茎は Milnor ファイバーのモノドロミーで、$s\equiv\alpha_j\bmod\Z$（$\alpha_j$ はスペクトル数）のとき 1 クラス、$s\in\Z$ のとき $H^0,H^1$ に寄与。全 7 種、$s\in\frac1{12}\Z$ または $\frac18\Z$、$|s|\le5$ で矛盾なし。Milnor モノドロミーによる説明は整合性の議論であり、証明ではない。$s=1$（$K_X$ のねじり）は茎が無限次元（$d_Q=0$）だが、大域的には $E_1$ 項の議論で扱える（命題 \ref{prop:K}）。
\end{obs}
\subsection{3重点の異常と大域予想}
\tV\ $\mathbb{CP}^2$ で $xy(x-y)$（3重点）：$m=1,2,4,5,\ldots$ と $m=-1$ は $[0,1,4,0,0]$、$m=-5,-7,-8$ は $[0,0,4,1,0]$、$m=-2,-3,-4$ は $[0,0,3,0,0]$。上の経験則で、$D_4$ のスペクトル数 $\{2/3,1,1,4/3\}$、$\mathrm{lct}=2/3$、$s=-m/3$：$m\ge-1$ は $s<2/3$ で $H^*(U;L)=(0,1,0)$ を加え $[0,1,4,0,0]$、$m\le-5$ は $s>2/3$ で $H_c^*=(0,0,0,1,0)$ を加え $[0,0,4,1,0]$、異常な $m=-2,-4$ は $s=2/3,4/3$ で、$\C^4$ の 1 クラス（次数 2）と $H_c$ の 1 クラス（次数 3）が差し引かれる（スペクトル系列 $H^p(X,\mathcal H^q)\Rightarrow\HH$ の $d_3$）。\tK\ この説明は予想である。尖点（$5/6,7/6$）、tacnode（$3/4,5/4$）、$D_5$（$5/8,\ldots$）のスペクトル数は $\frac13\Z$ に入らず、$\mathbb{CP}^1\times\mathbb{CP}^1$（$s\in\frac12\Z$）では $A_k$ のスペクトル数に $1$ 以外の半整数がなく、異常は現れない。

\tV\ 留数が等しくない場合：3重点で加群 $\prod h_i^{s_i}$（$h=zw(z-w)$）の茎は、533 個のランダムな $(s_1,s_2,s_3)\in(\frac16\Z\cap[-2,2])^3$（$s_i\ne1$）で $\C^4[-2]\oplus R$、$R=H^*(U_{\rm loc};L_s)$（$s_1+s_2+s_3<2$ のとき）、$0$（そうでないとき）。二重点では枝ごとの閾値であるのに対し、3重点の閾値は留数の\emph{和}である。\tK\ 一般原理：閾値は germ の対称トーラス作用に関する加群の重みと対数標準重みの比。

\begin{conj}[精密化 Riemann--Roch（等留数・孤立特異点）]\tK
$X$ コンパクト曲面、$\sigma\in H^0(-K_X)$ の零曲線 $D$ は簡約で孤立特異点のみ、$L$ は $c_1(L)=sK_X$ の Poisson 加群、$L_s$ は $U$ 上の階数 1 局所系（各枝まわりのモノドロミー $e^{2\pi is}$）、$\mathcal F_s$ は $U$ 上で $L_s$、$p\in D$ で $s<\mathrm{lct}_p$ なら $Rj_*L_s$ の茎、$s>\mathrm{lct}_p$ なら 0 とする構成可能複体。(a) $s\notin\mathrm{Sp}(D)$（$D$ の全特異点のスペクトル数と 1）のとき $\HH^*=H^*(X;\mathcal F_s)\oplus\bigoplus_p\C^{\mu_p}[-2]$。(b) Euler 標数は常に $\chi(X)$（$\chi(D)=\sum\mu_p$、$\delta$-不変量から\tP）。(c) $s\ll0$ では $\mathcal F_s=Rj_*L_s$ で、$\HH^*=H^*(U;L_s)\oplus\bigoplus\C^{\mu_p}[-2]$ は $s\bmod\Z$ のみに依存し次数 $\ge3$ で消える。(d) $s\in\mathrm{Sp}(D)$ では階数 1 のキャンセルが起こりうる。
\tV\ $\mathbb{CP}^2$（滑らか・結節・尖点・円錐＋割線・円錐＋接線・三角形・3重点）と $\mathbb{CP}^1\times\mathbb{CP}^1$（4 例）で、$m=-15..15$（$-9..8$）に対し (d) 以外で矛盾なし。証明されていない：茎の結果は準斉次 germ のみ、大域同定 $\mathbb H(C_m)\simeq(\mathcal F_s+\text{skyscraper})$ は当てはめ。非簡約 $D$（二重線、二重円錐）では茎は無限次元で、大域次元は $m$ について線形に増える（傾きは $\deg\mathcal O(1)|_C$、$\mathbb{CP}^2$ の二重線で 1、$\mathbb{CP}^1\times\mathbb{CP}^1$ の二重円錐で 2）。
\end{conj}

%--------------------------------------------------------------
\section{先行研究との比較}\label{sec:compare}
%--------------------------------------------------------------
\subsection{Hodge 理論（形式係数）}
Gualtieri は一般化 Kähler 多様体上の一般化 Hodge star、Born--Infeld 内積、Laplacian の等式、Hodge 分解、$dd^{\Jc}$-補題を示した（\tC\ Gualtieri, \emph{Generalized geometry and the Hodge decomposition}, arXiv:math/0409093；Gualtieri, \emph{Generalized K\"ahler geometry}, Comm.\ Math.\ Phys.\ 331 (2014) 297--331, arXiv:1007.3485；Cavalcanti, \emph{New aspects of the $dd^c$-lemma}, arXiv:math/0501406、および J.\ Geom.\ Phys.\ 57 (2006) 121--132。定理番号は精読していないので省く\tQ）。これらは自明束または $H$-捩れ係数の理論であり、直線束（曲率つき）係数の Bochner--Kodaira--Nakano 恒等式は、調べた範囲では書かれていない。本ノートの結果は、この係数付き版の穴を、Hodge 理論に頼らず、対数複体と構成可能層の比較で $\beta$-変換の場合に埋めたものと位置づけられる。

\subsection{Pal--Poddar}
\tC\ Pal--Poddar, \emph{Strong generalized holomorphic principal bundles}, arXiv:2404.18113 は、正則型（型 $k$ が一定）のコンパクト一般化複素多様体で、シンプレクティック葉構造の葉空間が Kähler 軌道体である場合に、強一般化正則（SGH）束の Hodge 理論、横断 Dolbeault 同型（Thm.\ 8.9, Cor.\ 8.10。番号は\tQ）、Hodge 分解と Serre 双対性（Thm.\ 11.2, 11.4, 11.6, 11.7）、Kähler の恒等式（Prop.\ 11.8, Lemma 11.9。番号は\tQ）、正の SGH 直線束（Def.\ 11.10, 11.11）に対する Kodaira 型・Serre 型の消滅定理（Thm.\ 11.12）を与える。正則型と軌道体の仮定は、葉構造が滑らかな葉束をなすことと、横断方向の Hodge 理論（楕円性）を保証するために使われている。本ノートとの違い：
\begin{itemize}
\item 型が跳ぶ場合（$\sigma$ の退化因子上で型が跳ぶ）は Pal--Poddar の枠外であり、本ノートの $\mathbb{CP}^2$ などの例はこの場合である。
\item SGH 束は一般化正則束の真部分クラス（葉に沿って平坦で横断方向に正則）。本ノートの Poisson 加群は $\beta$-変換構造における一般の GH 束である。
\item 定理 11.12 の主張・証明の構造は \S11 を再読して確認した。正確な主張：$H^q(M,(\mathcal G^*M)^p\otimes E)=0$ $(p+q>k)$、および $H^q(M,E'\otimes E^m)=0$ $(m\ge m_0,\ q>0)$。仮定（コンパクト、葉空間が Kähler 軌道体、$E$ が正）は、横断方向の Hodge 理論と Bochner--Kodaira--Nakano 型の曲率項の正値性に使われる。
\end{itemize}
原典の記述について、本調査で気づいた疑義（原典の誤植・省略の可能性が高いが未検証\tQ）：(i) 主張は $(i/2\pi)$ 付き、証明では $1/2\pi$ なし（無害）、(ii) $gc_1$ の符号の食い違いの原因は、特性多項式 (9.9) と \S10.1 の $\det(I\pm tA/2\pi i)$ の定義の食い違い、(iii)「正のクラスから正曲率の計量が出る」ことは Def.\ 11.11 の定義から証明されず、Thm.\ 11.12 の証明が仮定している、(iv) Def.\ 11.11 にコンパクト性・Kähler 性の明記がない、(v) Lemma 11.9(1) の証明は点ごとの局所枠での計算、(vi) Thm.\ 13.4(2)(3) の $T^{1,0}$ 対 $T^{1,0*}$、Example 13.10 の $q=n$ の二項係数。

\subsection{Tengzhou Hu の仕事}
\tC\ T.\ Hu, \emph{Kodaira vanishing theorems for K\"ahler Lie algebroids}, arXiv:2403.10876；\\ T.\ Hu, \emph{Lie groupoid Riemann--Roch--Hirzebruch theorem and applications}, arXiv:2403.13424。
\begin{itemize}
\item \emph{設定}（2403.13424）：Lie 亜群 $G\rightrightarrows M$（Def.\ 1、固有性・$s$-連結性の仮定は述べられていない）、その Lie 亜群 $\mathfrak g$。「複素 Lie 亜群」（Def.\ 2）は、\emph{実}の Lie 亜群 $\mathfrak g$ に、ファイバー方向の可積分な複素構造 $J$ を与えたもの（これは $\bT_\C$ の部分束としての複素 Lie 亜群 $L$ とは別物）。Kähler Lie 亜群は Def.\ 11、正の束は Def.\ 12。
\item \emph{主張}（2403.13424）：Thm.\ 4（Riemann--Roch）は、$M$ コンパクト、$\dbar_{\mathfrak g}$ が楕円型（主張されているが証明は見当たらない）、$\Omega$ が $\mathfrak g$-不変な密度（単模性）のとき
\[\mathrm{Ind}_\Omega(\dbar_{\mathfrak g}+\dbar_{\mathfrak g}^*)=(2\pi i)^{-k}\int_M(-1)^{\mathrm{rk}\,\mathfrak g}\bigl\langle\ch^{\mathfrak g}(\mathfrak g^{*p,0})\ch^{\mathfrak g}(E)\Td^{\mathfrak g}_\C(\mathfrak g),\Omega_{\mathfrak g}\bigr\rangle\]
を主張する。Thm.\ 7（消滅）：$E$ が正で $p+q>n$ なら $\mathfrak g$-Dolbeault Laplacian の\emph{核}が消える（コホモロジーの消滅は導けない：$\mathfrak g$ が一般の Lie 亜群のとき Laplacian は楕円型とは限らず、Hodge 理論が使えないため）。2403.10876 では Kähler Lie 亜群は Def.\ 2、Kähler の恒等式は Thm.\ 1 と Prop.\ 4、BKN 恒等式は Lemma 4、主な消滅定理は序論。
\item \emph{一般化 Dolbeault 複体は適用範囲外}\tP\ （導出、原典との照合は以下の範囲）：$L$ は複素 Lie 亜群で $J$ をもたず、$\wedge L^*$ は一般には $(p,q)$ 二重次数をもたない。正則 Poisson 構造 $\sigma$ に対し、実 Lie 亜群 $A_{\mathbb R}$ は Hu の設定に入るが、その Dolbeault 複体は $\bar\sigma$ から作られる $\wedge\mathfrak g^{*0,q}$ で、本ノートの複体 $(T^{0,1}\bowtie A^{1,0})$ とは別物。
\item \emph{指数は $m$ を見るか}\tK：Hu の指数は群の指数で、おおよそ $\int\ch(E)\Td(T^*)\ch(\wedge^pT)$ のように $m$ に多項式的に依存する見込みだが、本ノートの全 Euler 標数は $m$ によらない（命題 \ref{prop:euler}）。また周期的な跳びを決める留数 $-m/3$ は Hu のデータに含まれない。したがって Hu の定理は「一般化 Dolbeault コホモロジーに対する Riemann--Roch」ではない。
\item \emph{原典の不整合}：2403.13424（要約を通して得たもの）：Thm.\ 7 で $n$ が未定義、符号定数が Thm.\ 5, Prop.\ 1, Thm.\ 6 で異なる、「単模」の定義が二通り。2403.10876（本文の続きと付録の一部を読んで確認\tC）：Lemma 5 は $[L,\Lambda]=(p+q-n)\mathrm{id}$ だが、Thm.\ 2 の証明は $\Delta_{\bar\partial_{\mathfrak g}}=\Delta^{1,0}-[L,\Lambda]=\Delta^{1,0}+(p+q-n)\mathrm{id}$ と書き、後の等号の符号が合わない（$-(p+q-n)=n-p-q$）。さらに代入 $\nabla^2=\sqrt{-1}L$ は $\sqrt{-1}R=-L$ を意味し、「正」（$\sqrt{-1}R>0$、Def.\ 13）の定義と逆である。主張（正の $E$ で $p+q>n$ のとき核が消える）自体は古典の Kodaira--Nakano と整合する。つまり印刷された証明に符号のずれがある（主張の否定ではない）。Cor.\ 1 に矢印と記号 $s/t$ の誤植。付録（Lemma 6--10、$[\bar\partial_{\mathfrak g},\iota_\omega]$ の計算）は添字が読み取りにくく検証できていない。参考文献欄と掲載誌は未確認。\tQ
\item \emph{本ノートの前版（予備調査）の訂正}：前版では著者を「S.\ Hu」、題名を別の題名で書いていたが、いずれも誤りで、上記が正しい。
\end{itemize}

\subsection{積分可能性・群}
\tC\ Gualtieri--Li（arXiv:1206.3674、IMRN 2014）は対数シンプレクティック多様体の Hausdorff シンプレクティック亜群を分類し、滑らかな $D$ に沿って接する対数ベクトル場の Lie 亜群 $T_X(-\log D)$ が積分可能であることを与える。Bailey--Gualtieri（arXiv:1611.03850）は、一般化複素構造の弱正則シンプレクティック亜群への積分と実 Poisson 構造の積分可能性の同値を示す（Thm.\ 7.1：下の実 Poisson 構造が積分可能であることが必要十分。Thm.\ 7.1 自体は Hausdorff・$s$-連結性を述べない。構成 Thm.\ 7.6 は $s$-連結かつ $s$-単連結な積分を用いる）。Laurent-Gengoux--Stiénon--Xu（arXiv:0803.2031）：Thm.\ 3.17, 3.22。Crainic--Fernandes（arXiv:math/0105033）：Thm.\ 4.1, Cor.\ 5.2。これらは、指数理論の入力としての群を与えるが、複素 Lie 亜群 $L$ 自体の積分は扱わない。

%--------------------------------------------------------------
\section{未解決問題と次のステップ}\label{sec:open}
%--------------------------------------------------------------
\begin{enumerate}
\item \textbf{三種の神器のうち未着手のもの}：Kodaira 埋め込みの一般化。$\HH^0$ が Casimir 切断であること（観察 \ref{obs:casimir}）から、$M$ 自身への埋め込みではなく、Casimir 切断が分離する商（葉空間など）への写像として定式化する必要がある。Casimir 切断の線形系のベースポイント・分離性を、$\mathbb{CP}^2$ の例（$\HH^0=\C F^{m/3}$、すなわち $m\in3\Z_{\ge0}$ で 1 次元）から調べるのが次の具体的な一歩。
\item \textbf{Serre 型消滅の一般化}（予想 \ref{conj:A3}）：「正」な一般化正則直線束の定義、$\beta$-変換でない $\Jc$ での $E_1$ 項に相当する濾過の構成。Pal--Poddar Thm.\ 11.12 との比較は \S11 の再読で行った（群は別物、横断次数の読みは整合）。
\item \textbf{Riemann--Roch の定式化}：全 Euler 標数ではなく、細分（$\Jc$ の型による濾過の $E_1$ 項、周期擬多項式、指標根 $H^*(\partial U)$）を決める。予想 \ref{conj:I3} の証明には、複素 Lie 亜群 $L$ の指数理論（単模性・Fredholm 性 $\Leftrightarrow r\notin\Z$ の導出）が必要。境界つき一般化複素多様体への APS 型の拡張は、さらに先の課題。
\item \textbf{対数シンプレクティック以外}：非ノーダル特異点の局所理論の証明（準斉次 germ のみ計算）、留数が等しくない場合の大域公式、非簡約 $D$（無限次元の茎と線形増大）、次元 4 以上の非トーリックな例（観察 \ref{obs:arith} を参照）。
\item \textbf{一般化 Kähler 側}：直線束係数の Bochner--Kodaira--Nakano 恒等式を一般化 Kähler の枠組みで書き下す。$H$-捩れと曲率項の相互作用が障害になりうる\tK。
\item \textbf{本ノートの未確認事項の解消}：\S\ref{sec:ledger} の \tQ\ の項目。特に、Mokri の定理、LGSX Lemma 4.22 の再導出、Gualtieri 2011 \S4 の精読、残る書誌情報。
\end{enumerate}

\paragraph{次のステップの優先順位（筆者案）。}
(a) 型が跳ばない比較対象として、Pal--Poddar の正則型の例で同じ「Poincaré 多項式」を計算して、本ノートの周期的な跳びと見比べる。(b) $\mathbb{CP}^2$ の例について、$\HH^0=\C F^{m/3}$ を使った Casimir 切断の線形系の幾何を調べ、Kodaira 型定理の「正しい形」を探る。(c) 予想 \ref{conj:I3} の指数理論的な証明の足場として、実 Lie 亜群 $T_X(-\log D)$ 上の重みつき $b$-複体と留数 $r$ の関係を厳密に定式化する。

%--------------------------------------------------------------
\section{主張の状態一覧と検証の範囲}\label{sec:ledger}
%--------------------------------------------------------------
\subsection{主張ごとの状態}
\begin{center}\footnotesize
\begin{tabular}{>{\raggedright\arraybackslash}p{4.0cm}>{\raggedright\arraybackslash}p{3.0cm}>{\raggedright\arraybackslash}p{8.4cm}}\toprule
主張 & 状態 & 依存・穴\\\midrule
定理 \ref{thm:main1}（Poisson 加群 $=$ 対数接続） & \tP & 任意の偶数次元・単純正規交差。留数公式 $c_1=-\sum r_i[D_i]$ は Esnault--Viehweg の引用。\\
定理 \ref{thm:cp2}（$\mathbb{CP}^2$ の表） & \tP\ \tV & 手計算の導出と $m=-15..15$ の厳密計算が一致。複体の同一視（命題 \ref{prop:ident}(b)）に依存。\\
定理 \ref{thm:main2}（曲面の主定理） & \tP & 命題 \ref{prop:ident}(b)（LGSX Thm.\ 4.19 と Mokri の定理）、Deligne の比較、局所コホモロジーの Laurent 尾の計算（標準だが未記述）。$r_i\ne1$ が必要。\\
定理 \ref{thm:B}（閉公式） & \tP\ \tV & 定理 \ref{thm:main2} に加え (H1)--(H4)。(H4)（$H^1(U;L_\chi)=0$）は 8 例で確認、一般には $V^1(U)$ の条件。196 個の値で一致。\\
定理 \ref{thm:A1}（Serre 型） & \tP & 古典 Serre（Hartshorne III.5.2）と非輪状分解のみ。任意次元。\\
定理 \ref{thm:A2}（周期的安定化） & \tP & 定理 \ref{thm:main2} とその仮定（$X$ 単連結、$U$ アファイン）に依存。\\
命題 \ref{prop:euler}（Euler 標数） & \tP & HRR のみ。\\
定理 \ref{thm:dual}（双対性） & \tP\ \tV & 補題 D, E は $n\le4$ で機械確認。$\dbar$ 部分と $d$ 部分の Koszul 符号の整合性は手計算のみ。楕円型複体の Hodge 理論は標準。\\
定理 \ref{thm:weight}（weight 公式） & \tP\ \tV & $E_1^{a>0}=0$ を仮定。60 例の線形代数との一致、3626 例の一般 $A$ での一致。\\
観察 \ref{obs:arith}（$\sigma$ の算術依存） & \tV & 4 次元の少数の例。機構の説明は概略。\\
\S\ref{sec:local} の局所理論・3 重点の説明 & \tV\ \tK & 準斉次 germ のみ。大域同定は当てはめ。\\
予想 \ref{conj:A3}, \ref{conj:I3} & \tK & 定義・証明ともに未了。\\
Kodaira 埋め込みの一般化 & --- & 扱っていない（\S\ref{sec:open}）。\\\bottomrule
\end{tabular}
\end{center}

\subsection{未確認・要注意の項目（\tQ）}
\begin{itemize}
\item \textbf{書誌情報（記憶に基づく）}：Goto の掲載誌・巻頁、Gualtieri--Li の IMRN 巻頁、Gualtieri ``Branes'' の OUP 書誌、Mokri の巻頁、Arapura の巻頁、Lang--Jia--Liu の掲載誌、LGSX の IMRN Art.\ ID、Pal--Poddar の掲載巻・頁。Deligne LNM 163、Crainic--Fernandes、Fulton、Esnault--Viehweg の番号。Pal--Poddar の Thm.\ 8.9, Prop.\ 9.3, Ex.\ 10.18, Prop.\ 11.8, Lemma 11.9, Ex.\ 13.10 の番号と内容。（v2 で確認済み：Hitchin Moscow Math.\ J.\ 12 (2012) 567--591、Pym Lett.\ Math.\ Phys.\ 108 (2018) 573--632、Hausel--Hunsicker--Mazzeo Duke 122 (2004) 485--548、Timmerscheidt J.\ reine angew.\ Math.\ 379 (1987) 152--171、Pal--Poddar の題名・著者。）
\item \textbf{Gualtieri 2011（arXiv:math/0703298）}：約 16 万字のうち先頭約 10 万字（\S3 の途中まで）を精読。\S4（Darboux の定理など）は未精読。v2 で、型 Def.\ 3.5、標準束 Def.\ 3.7、楕円性 Prop.\ 3.13/Cor.\ 3.14、Picard 群 Def.\ 3.18、Darboux Thm.\ 4.3（正則点）、Def.\ 4.1（正則点の定義）を確認した。型の跳びの例は \S4.1、$\rho=z_1+dz_1\wedge dz_2$（$\R^4$）。
\item \textbf{Hitchin, \emph{Poisson modules and generalized geometry}（arXiv:0905.3227、題名・著者は確認済み）}：GH 束 $=$ Poisson 加群を序論で述べるのみで、$\Jc_\sigma$ の記号、超コホモロジーの同定、Atiyah 類による障害、$\mathbb{CP}^2$ の例、Lie 亜群 $L$ は含まれない（取得した本文による）。本ノートは代わりに Branes on Poisson varieties（arXiv:0710.2719）と LGSX を引用した。（v1 は著者を Gualtieri と誤記し、「確認した」としていたが、著者名と内容の両面で誤りだった。）
\item \textbf{Pic$(\Jc)=\mathbb H^1(\mathcal L_{\log})$}：どの文献でも本文を確認できなかった。本ノートは使っていない。使ったのは、Poisson 加群構造の存在と一意性の直接導出のみ。
\item \textbf{LGSX Lemma 4.22、Mokri の定理}：再導出していない。Hodge 理論・楕円型複体の標準的事実と、Esnault--Viehweg、Deligne、Fulton--Deligne、Timmerscheidt、Arapura、Crainic--Fernandes は引用であり、適用範囲を原典で確認していない。
\item \textbf{Hu の 2 論文}：2403.10876 は先頭約 20 万字（全体は約 50 万字）まで、2403.13424 は概要と定理の主張のみ。
\item \textbf{Pal--Poddar}：\S\S5--11, 13 を再読した（2026-11-22）。抽出は要約を介しており、式の細部は未確認。
\item \textbf{計算の限界}：厳密な線形代数は有限の範囲の $m$、有限個の曲線（Fermat 型・Hesse 型の 3 次曲線など）、素数を法とする計算（$32003$、$2^{61}-1$）。曲線の取り方への依存（ランダムな滑らかな曲線）と、法 $p$ での還元が標数 0 の結果を代表することは検証していない。
\end{itemize}

\subsection{作業の記録}
スクリプトはすべて Python（標準ライブラリ）による厳密算術で、\texttt{cp2/} に保存した（\texttt{poisson.py, table.py, neg\_check.py, predict.py, sing\_table.py, p1p1.py, local\_stalk.py, local\_unequal.py, ncd\_cochain\_check.py, ncd\_local\_weights.py, ncd\_global\_check.py, serre\_rr\_check.py, serre\_dual\_check.py, duality\_ndim.py, duality\_ndim\_exact.py, toric\_hh.py, toric\_crit.py, toric\_generic\_check2.py, index\_route\_check.py}）。

%--------------------------------------------------------------
\section{v2 改訂点（v1 $\to$ v2、2026年11月22日）}\label{sec:rev}
%--------------------------------------------------------------
v1 の引用の再確認（arXiv ページ・出版社ページ・原文の再読）で見つかった誤りと更新点。数学的な主張（定理・命題）は変更していない。
\begin{enumerate}
\item \textbf{arXiv:0905.3227 の著者}：Gualtieri ではなく \textbf{Hitchin}（\emph{Poisson modules and generalized geometry}）。v1 の参考文献と \S\ref{sec:ledger} の記述を訂正。
\item \textbf{Gualtieri 2011 の番号}：楕円性は Prop.\ 3.13 と Cor.\ 3.14（v1 の「Prop.\ 3.16」は存在しない）。Def.\ 3.18 は Picard 群の定義で、一般化正則束は Thm.\ 3.15 の後に散文で導入される。一般化 Darboux は Thm.\ 4.3（v1 の「4.35」は誤り）。v1 の「例 4.1--4.2」は存在せず、Def.\ 4.1 が正則点の定義。
\item \textbf{Pal--Poddar（arXiv:2404.18113）}：題名 \emph{Strong generalized holomorphic principal bundles}、著者 D.\ Pal と M.\ Poddar。v1 の「著者表示の食い違い」は再現しなかった。定理番号を整理（Hodge 分解・Serre 双対性は Thm.\ 11.2, 11.4, 11.6, 11.7、消滅は Thm.\ 11.12）。Thm.\ 11.12 の主張と、「消滅は横断次数についてのみ」という本ノートの読みとの整合を、\S\S5--11, 13 の再読で確認（群は同一ではない）。原文の疑義の原因（特性多項式の定義の食い違い）を更新。
\item \textbf{Lang--Jia--Liu}：arXiv:2102.07177、\emph{The Atiyah class of generalized holomorphic vector bundles}。局所自明化で定義される GH 束の論文で、Pal--Poddar の SGH 束とは別。
\item \textbf{Hu（Tengzhou Hu）2403.10876}：印刷された消滅定理の証明に符号のずれがあることを本文の再読で確認した（主張自体は古典と整合）。
\item \textbf{書誌の確定}：Hitchin (Moscow Math.\ J.\ 12 (2012) 567--591)、Pym (Lett.\ Math.\ Phys.\ 108 (2018) 573--632)、Hausel--Hunsicker--Mazzeo (Duke 122 (2004) 485--548)、Timmerscheidt (J.\ reine angew.\ Math.\ 379 (1987) 152--171、v1 の記憶は誤り)、Mokri の年・DOI、Goto の arXiv 番号 1501.03398、Arapura の arXiv 番号。
\item \textbf{Bailey--Gualtieri Thm.\ 7.1}：実 Poisson 構造の積分可能性が必要十分であることを確認し、Hausdorff・$s$-連結性に関する注意を追記。
\end{enumerate}
なお、本節で確認済みとした以外の \tQ\ 項目は、\S\ref{sec:ledger} に残している。

\begin{thebibliography}{99}\small
\bibitem{AB} M.\ Abouzaid, M.\ Boyarchenko, Local structure of generalized complex manifolds, J.\ Symplectic Geom.\ 4 (2006) 43--62, arXiv:math/0511049.
\bibitem{AN} Y.\ Akizuki, S.\ Nakano, Note on Kodaira--Spencer's proof of Lefschetz's theorems, Proc.\ Japan Acad.\ 30 (1954) 266--272.
\bibitem{ALN} B.\ Ammann, R.\ Lauter, V.\ Nistor, Pseudodifferential operators on manifolds with a Lie structure at infinity, arXiv:math/0304044.
\bibitem{Ara} D.\ Arapura, Geometry of cohomology support loci for local systems I, arXiv:alg-geom/9612006; J.\ Algebraic Geom.\ 6 (1997) 563--597 [巻頁は記憶\tQ].
\bibitem{AS} M.\ Atiyah, I.\ Singer, The index of elliptic operators III, Ann.\ Math.\ 87 (1968) 546--604.
\bibitem{BG} A.\ Bailey, M.\ Gualtieri, Integration of generalized complex structures, arXiv:1611.03850.
\bibitem{BLS} M.\ Braverman, Y.\ Loizides, Y.\ Song, Geometric quantization of $b$-symplectic manifolds, arXiv:1910.10016.
\bibitem{Cav} G.\ Cavalcanti, New aspects of the $dd^c$-lemma, arXiv:math/0501406; The decomposition of forms and cohomology of generalized complex manifolds, J.\ Geom.\ Phys.\ 57 (2006) 121--132.
\bibitem{CF} M.\ Crainic, R.\ L.\ Fernandes, Integrability of Lie brackets, Ann.\ Math.\ 157 (2003) 575--620, arXiv:math/0105033.
\bibitem{Del} P.\ Deligne, \'Equations diff\'erentielles \`a points singuliers r\'eguliers, Lecture Notes in Math.\ 163, Springer 1970.
\bibitem{DelFG} P.\ Deligne, Le groupe fondamental du compl\'ement d'une courbe plane n'ayant que des points doubles ordinaires est ab\'elien, S\'em.\ Bourbaki 543 (1979/80).
\bibitem{EV} H.\ Esnault, E.\ Viehweg, Logarithmic de Rham complexes and vanishing theorems, Invent.\ Math.\ 86 (1986) 161--194.
\bibitem{Ful} W.\ Fulton, On the fundamental group of the complement of a node curve, Ann.\ Math.\ 111 (1980) 407--409.
\bibitem{Goto} R.\ Goto, Unobstructed deformations of generalized complex structures induced by $C^\infty$ logarithmic symplectic structures, arXiv:1501.03398 [掲載誌・巻頁は未確認\tQ].
\bibitem{GH} P.\ Griffiths, J.\ Harris, Principles of Algebraic Geometry, Wiley 1978.
\bibitem{GuaTh} M.\ Gualtieri, Generalized complex geometry, Oxford DPhil thesis, arXiv:math/0401221.
\bibitem{GuaHodge} M.\ Gualtieri, Generalized geometry and the Hodge decomposition, arXiv:math/0409093.
\bibitem{Gua11} M.\ Gualtieri, Generalized complex geometry, Ann.\ Math.\ 174 (2011) 75--123, arXiv:math/0703298.
\bibitem{GuaK} M.\ Gualtieri, Generalized K\"ahler geometry, Comm.\ Math.\ Phys.\ 331 (2014) 297--331, arXiv:1007.3485.
\bibitem{GuaBr} M.\ Gualtieri, Branes on Poisson varieties, arXiv:0710.2719 (v2, 2010-07；arXiv ページに掲載書誌の記載なし。The Many Facets of Geometry, Oxford Univ.\ Press 2010 は記憶\tQ).
\bibitem{GuaPM} N.\ Hitchin, Poisson modules and generalized geometry, arXiv:0905.3227 (2009)（題名・著者は arXiv ページで確認。掲載書誌は未確認\tQ）.
\bibitem{GL} M.\ Gualtieri, S.\ Li, Symplectic groupoids of log symplectic manifolds, IMRN 2014, arXiv:1206.3674.
\bibitem{Hart} R.\ Hartshorne, Algebraic Geometry, GTM 52, Springer 1977.
\bibitem{HHM} T.\ Hausel, E.\ Hunsicker, R.\ Mazzeo, Hodge cohomology of gravitational instantons, Duke Math.\ J.\ 122 (2004) 485--548, arXiv:math/0207169.
\bibitem{Hirz} F.\ Hirzebruch, Neue topologische Methoden in der algebraischen Geometrie, Springer 1956.
\bibitem{Hit03} N.\ Hitchin, Generalized Calabi--Yau manifolds, Q.\ J.\ Math.\ 54 (2003) 281--308, arXiv:math/0209099.
\bibitem{Hit12} N.\ Hitchin, Deformations of holomorphic Poisson manifolds, Moscow Math.\ J.\ 12 (2012), no.\ 3, 567--591, arXiv:1105.4775.
\bibitem{Hu1} T.\ Hu, Kodaira vanishing theorems for K\"ahler Lie algebroids, arXiv:2403.10876.
\bibitem{Hu2} T.\ Hu, Lie groupoid Riemann--Roch--Hirzebruch theorem and applications, arXiv:2403.13424.
\bibitem{Huy} D.\ Huybrechts, Complex Geometry: An Introduction, Springer 2005.
\bibitem{Kod53} K.\ Kodaira, On a differential-geometric method in the theory of analytic stacks, Proc.\ Nat.\ Acad.\ Sci.\ USA 39 (1953) 1268--1273.
\bibitem{Kod54} K.\ Kodaira, On K\"ahler varieties of restricted type, Ann.\ Math.\ 60 (1954) 28--48.
\bibitem{Lan} M.\ Lanius, Poisson cohomology of a class of log symplectic manifolds, arXiv:1605.03854.
\bibitem{LGSX1} C.\ Laurent-Gengoux, M.\ Sti\'enon, P.\ Xu, Holomorphic Poisson manifolds and holomorphic Lie algebroids, IMRN 2008, arXiv:0707.4253.
\bibitem{LGSX2} C.\ Laurent-Gengoux, M.\ Sti\'enon, P.\ Xu, Integration of holomorphic Lie algebroids, arXiv:0803.2031.
\bibitem{LJL} H.\ Lang, X.\ Jia, Z.\ J.\ Liu, The Atiyah class of generalized holomorphic vector bundles, arXiv:2102.07177; Differential Geom.\ Appl.\ 90 (2023) 102031（掲載誌情報は Pal--Poddar の参考文献による\tQ）。局所自明化と一般化正則な遷移関数で定義される GH 束の Atiyah 類を扱う。Pal--Poddar の SGH 束とは別の文献。
\bibitem{MM} R.\ Mazzeo, R.\ Melrose, Pseudodifferential operators on manifolds with fibred boundaries, arXiv:math/9812120.
\bibitem{Mel} R.\ Melrose, The Atiyah--Patodi--Singer Index Theorem, A K Peters 1993.
\bibitem{Mok} T.\ Mokri, Matched pairs of Lie algebroids, Glasgow Math.\ J.\ 39 (1997) 167--181, doi:10.1017/S0017089500032055 [年・DOI は出版社ページで確認、巻頁は記憶\tQ].
\bibitem{PP} D.\ Pal, M.\ Poddar, Strong generalized holomorphic principal bundles, arXiv:2404.18113; Ann.\ Sc.\ Norm.\ Super.\ Pisa Cl.\ Sci., doi:10.2422/2036-2145.202501\_026（掲載巻号は\tQ）.
\bibitem{PPT} M.\ Pflaum, H.\ Posthuma, X.\ Tang, The index of geometric operators on Lie groupoids, Indag.\ Math.\ 25 (2014) 1--18, arXiv:1308.0236.
\bibitem{Pym} B.\ Pym, Constructions and classifications of projective Poisson varieties, Lett.\ Math.\ Phys.\ 108 (2018), no.\ 3, 573--632, doi:10.1007/s11005-017-0984-5.
\bibitem{QS} Y.\ Qiao, Y.\ So, Deformation and K-theoretic index formulae on boundary groupoids, arXiv:2304.13345.
\bibitem{Serre} J.-P.\ Serre, Faisceaux alg\'ebriques coh\'erents, Ann.\ Math.\ 61 (1955) 197--278.
\bibitem{Tim} K.\ Timmerscheidt, Mixed Hodge theory for unitary local systems, J.\ reine angew.\ Math.\ 379 (1987) 152--171.
\end{thebibliography}


\end{document}
